% !TEX program = xelatex
\documentclass[11pt,a4paper]{article}

% Encoding and fonts
\usepackage{fontspec}
\usepackage{xeCJK}
\setCJKmainfont[Path=fonts/noto-cjk/]{NotoSansCJKsc-Regular.otf}

% Math
\usepackage{amsmath,amssymb}

% Layout
\usepackage{geometry}
\geometry{tmargin=2cm,bmargin=2cm,lmargin=1.8cm,rmargin=1.8cm}

% Parallel columns for bilingual text
\usepackage{paracol}
\columnsep=1.5em
\globalcounter{section}
\globalcounter{subsection}

% Typography
\usepackage{titlesec}
\usepackage{setspace}
\usepackage{indentfirst}
\setlength{\parindent}{2em}
\linespread{1.25}

% Colors for labels
\usepackage{xcolor}
\definecolor{classicallabel}{RGB}{128,0,0}
\definecolor{modernlabel}{RGB}{0,64,128}
\definecolor{volumecolor}{RGB}{40,40,100}

% Custom commands
\newcommand{\Volume}[1]{%
  \switchcolumn[0]*
  \vspace{1.5em}
  \begin{center}
    {\color{volumecolor}\rule{0.6\columnwidth}{0.6pt}}\\[0.8em]
    {\Large\bfseries\color{volumecolor} 卷#1}\\[0.3em]
    {\color{volumecolor}\rule{0.6\columnwidth}{0.6pt}}
  \end{center}
  \vspace{0.8em}
  \switchcolumn[1]
  \vspace{1.5em}
  \begin{center}
    {\color{volumecolor}\rule{0.6\columnwidth}{0.6pt}}\\[0.8em]
    {\Large\bfseries\color{volumecolor} 第#1卷}\\[0.3em]
    {\color{volumecolor}\rule{0.6\columnwidth}{0.6pt}}
  \end{center}
  \vspace{0.8em}
}

\newcommand{\Section}[1]{%
  \switchcolumn[0]*
  \subsection*{#1}
  \switchcolumn[1]
  \subsection*{\mdseries\color{modernlabel} #1}
}

\newcommand{\Question}[2]{%
  \switchcolumn[0]*
  \noindent{\bfseries\color{classicallabel} #1}\\#2
  \switchcolumn[1]
  \noindent{\bfseries\color{modernlabel} 【问】}\\#2
}

\newcommand{\Shi}[2]{%
  \switchcolumn[0]*
  \noindent{\bfseries\color{classicallabel} #1}\\#2
  \switchcolumn[1]
  \noindent{\bfseries\color{modernlabel} 【释】}\\#2
}

\newcommand{\Shu}[2]{%
  \switchcolumn[0]*
  \noindent{\bfseries\color{classicallabel} #1}\\#2
  \switchcolumn[1]
  \noindent{\bfseries\color{modernlabel} 【术】}\\#2
}

\newcommand{\You}[2]{%
  \switchcolumn[0]*
  \noindent{\bfseries\color{classicallabel} #1}\\#2
  \switchcolumn[1]
  \noindent{\bfseries\color{modernlabel} 【又】}\\#2
}

\newcommand{\YouShu}[2]{%
  \switchcolumn[0]*
  \noindent{\bfseries\color{classicallabel} #1}\\#2
  \switchcolumn[1]
  \noindent{\bfseries\color{modernlabel} 【又术】}\\#2
}

\newcommand{\Note}[2]{%
  \switchcolumn[0]*
  \noindent{\itshape\small #1}\\#2
  \switchcolumn[1]
  \noindent{\itshape\small\color{gray} #1}\\#2
}

\newcommand{\EndVol}[1]{%
  \switchcolumn[0]*
  \vspace{1em}
  \begin{center}
    \rule{0.5\columnwidth}{0.4pt}\\[0.5em]
    {\bfseries #1}
  \end{center}
  \switchcolumn[1]
  \vspace{1em}
  \begin{center}
    \rule{0.5\columnwidth}{0.4pt}\\[0.5em]
    {\bfseries\color{gray} #1}
  \end{center}
}

% No page numbers for title page
\pagenumbering{gobble}

\begin{document}

% ======================== TITLE PAGE ========================
\begin{center}
\vspace*{2em}
{\Huge\bfseries 测圆海镜分类释术}\\[0.8em]
{\Large\color{volumecolor} 卷六至卷十}\\[1.5em]
{\large 文言文 \quad $\bullet$ \quad 白话文对照本}\\[2em]
{\large 元\quad 李\quad 冶\quad 撰}\\[0.3em]
{\large 明\quad 顾应祥\quad 释术}\\[3em]
{\color{volumecolor}\rule{0.5\textwidth}{1pt}}\\[1.5em]
{\normalsize 左栏：原文（文言文）\quad | \quad 右栏：今译（白话文）}
\end{center}

\vspace{2em}

\begin{minipage}{0.9\textwidth}
\setlength{\parindent}{2em}
《测圆海镜》乃金元之际数学家李冶（1192--1279）所著，为中国古代天元术之杰作。全书十二卷，以勾股容圆为题，设问一百七十道，系统阐述了以天元术（代数方法）求解几何问题之法。明顾应祥为之分类释术，编成《测圆海镜分类释术》。

本文档收录卷六至卷十（最终五卷）之内容。卷六论通弦测望，卷七论大差小差，卷八论皇极太虚，卷九论诸弦差较，卷十论杂法。
\end{minipage}

\vspace{2em}
\begin{center}
{\color{volumecolor}\rule{0.5\textwidth}{1pt}}
\end{center}

\newpage
\tableofcontents
\newpage

% ======================== START BILINGUAL COLUMNS ========================
\pagenumbering{arabic}
\setcounter{page}{1}

\begin{paracol}{2}

% ---- Column headers ----
\switchcolumn[0]
\begin{center}
{\large\bfseries\color{classicallabel} 原文（文言文）}
\end{center}
\vspace{0.5em}

\switchcolumn[1]
\begin{center}
{\large\bfseries\color{modernlabel} 白话译文}
\end{center}
\vspace{0.5em}

\switchcolumn[0]*
\noindent\rule{\columnwidth}{0.3pt}
\switchcolumn[1]
\noindent\rule{\columnwidth}{0.3pt}

%======================================================================
\Volume{六}
%======================================================================

\switchcolumn[0]*
\noindent\textbf{钦定四库全书}

\noindent\textbf{测圆海镜分类释术卷六}\quad 元~李~冶~撰\\
明~顾应祥~释术

\switchcolumn[1]
\noindent\textbf{钦定四库全书}

\noindent\textbf{测圆海镜分类释术·第六卷}\quad 元代~李~冶~撰\\
明代~顾应祥~释术

\vspace{1em}

%----------------------------------------------------------------------
\Section{勾与和测望一}
%----------------------------------------------------------------------

\Question{问一}{%
甲乙俱在城外西北乾隅，甲南行不知步数而立，乙东行三百二十步见之，甲又斜行与相会，计甲直行斜行共一千二百八十步。问城径。
}

\switchcolumn[0]*
\noindent\textbf{释曰}\quad 此通勾与通股弦和测望。乙东行，通勾也；甲直斜共行，通股弦和也。

\switchcolumn[1]
\noindent\textbf{【释】}\quad 这是用通勾与通股弦和来进行测望的问题。乙向东行走，是通勾；甲直行与斜行一共行走的步数，是通股弦和。

\switchcolumn[0]*
\noindent\textbf{术曰}\quad 勾自之，得一十万二千四百。以和除之，得八十为股弦较。以较减和，半之为股。以勾股求容圆术求之，得城径。

\switchcolumn[1]
\noindent\textbf{【术】}\quad 将勾自乘，得到十万二千四百。用这个数除以和，得到八十，这就是股弦差。用和减去这个差，取一半得到股。再用勾股求容圆的公式计算，求得城圆的直径。

\switchcolumn[0]*
\noindent\textbf{又曰}\quad 勾和各自乘，相减为实；倍和除之，得股。相并为实，倍和除之，得弦。

\switchcolumn[1]
\noindent\textbf{【又】}\quad 将勾与和分别自乘，相减作为被除数；用二倍和来除，得到股。将勾与和相加后自乘作为被除数，用二倍和来除，得到弦。

\switchcolumn[0]*
边勾以下，俱以类推。

\switchcolumn[1]
边勾以下的各类问题，都按照这个方法类推。

\vspace{0.5em}

\Question{问二}{%
乙出东门南行，丙出南门东行，各不知步数而立，只云丙行多于乙步。甲从乾隅东行三百二十步，望乙丙与城相参直，计乙丙共行一百二步。问城径。
}

\switchcolumn[0]*
\noindent\textbf{释曰}\quad 此以通勾与明勾、边股和测望。甲东行，通勾也；乙出东门南行为边股，丙出南门东行为明勾，共计一百二步，明勾、边股和也。

\switchcolumn[1]
\noindent\textbf{【释】}\quad 这是用通勾与明勾、边股和来进行测望的问题。甲向东行走，是通勾；乙出东门向南行走是边股，丙出南门向东行走是明勾，共计一百零二步，就是明勾与边股的和。

\switchcolumn[0]*
\noindent\textbf{术曰}\quad 倍共步乘东行算，得二千零八十八万九千六百，为立方实。共步乘东行，加东行算，得一十三万五千零四十，为从方。东行为从廉，五分为隅算。作带从负隅，以廉减从，开立方法除之，得全径。

\switchcolumn[1]
\noindent\textbf{【术】}\quad 将共步的二倍乘以向东行走的数，得到二千零八十八万九千六百，作为立方的被除数。共步乘以向东行走的数，加上东行数的平方，得到十三万五千零四十，作为从方。东行数作为从廉，五分作为隅算。按带从负隅的方法，用廉来减从方，用开立方法除之，得到圆的全径。

\switchcolumn[0]*
带从负隅，以廉减从，半翻法开立方曰：置所得实，以从方约之。初商二百，置一于左上为法，置一乘从廉得六万四千，以减从方，存七万一千零四十为从。置一自之得四万，以隅算五分因之，得二万为隅法。并从，共九万一千零四十为下法，与上法相乘，除实一千八百二十万八千，余实二百六十八万一千六百。从方内再减六万四千，止余七千零四十为从。三因隅法得六万为方法，三因初商得六百为廉法。次商四十，置一于左，次为上法，置一乘从廉得一万二千八百，以减余从，不及减，反减余从七千零四十，余五千七百六十为负从。置一乘廉法，以隅因得一万二千，置一自之，隅因得八百为隅法。并方廉隅，共七万二千八百，减去负从，余六万七千零四十为下法，与上法相乘，除实尽。

\switchcolumn[1]
\noindent\textbf{【带从负隅开立方详解】}\quad 把求得的被除数，用从方来约。初次商得二百，在左上方放一个一作为上法，用一乘以从廉得到六万四千，用来减从方，剩余七万一千零四十作为新的从方。一自乘得四万，用隅算五分乘之，得到二万作为隅法。合并从方与隅法，共九万一千零四十作为下法，与上法相乘，从被除数中减去一千八百二十万八千，剩余二百六十八万一千六百。从方中再减去六万四千，只剩下七千零四十作为从方。将隅法三倍得六万作为方法，将初商三倍得六百作为廉法。第二次商得四十，在左上方放一个一作为上法，用一乘以从廉得到一万二千八百，用来减剩余的从方，不够减，反过来用剩余的从方七千零四十去减，余五千七百六十作为负从。用一乘以廉法，再用隅算乘之得到一万二千，一自乘再用隅算乘之得到八百作为隅法。合并方法、廉法、隅法共七万二千八百，减去负从五千七百六十，余六万七千零四十作为下法，与上法相乘，把剩余被除数刚好除尽。

\switchcolumn[0]*
法已见四卷通勾太虚弦条，因以五分为隅，故重出。

\switchcolumn[1]
这个方法在第四卷通勾太虚弦条下已经出现过，因为这里用五分作为隅算，所以重新列出。

\switchcolumn[0]*
又为带从负隅，以廉添积，开立方法，法见四卷通勾虚弦条下。

\switchcolumn[1]
又可以用带从负隅的方法，以廉来增添积数，开立方法求解，详见第四卷通勾虚弦条下。

\vspace{0.5em}

\Question{问三}{%
乙出东门东行，丙出南门南行，各不知步数而立。甲从乾隅东行三百二十步，望乙丙二人俱与城相参直，计乙丙共行一百五十一步。问城径。
}

\switchcolumn[0]*
\noindent\textbf{释曰}\quad 此以通勾与边勾、明股和，立法测望。甲东行，通勾；乙东行，边勾；丙南行，明股也。

\switchcolumn[1]
\noindent\textbf{【释】}\quad 这是用通勾与边勾、明股和来建立方法进行测望的问题。甲向东行走，是通勾；乙向东行走，是边勾；丙向南行走，是明股。

\switchcolumn[0]*
\noindent\textbf{术曰}\quad 通勾自之，得一十万二千四百，半之得五万一千二百，又自之得二十六亿二千一百四十四万，为三乘方实。以三百六十二乘半通勾算，得一千八百五十三万四千四百，为从方。通勾乘和步，得四万八千三百二十，为从一廉。五之通勾，得一千六百，为从二廉。二分五厘为常法。作带从方廉三乘方法，开之得八十为小差。小差者，通股弦较也。以减通勾，即城径。

\switchcolumn[1]
\noindent\textbf{【术】}\quad 将通勾自乘，得到十万二千四百，取一半得到五万一千二百，再自乘得到二十六亿二千一百四十四万，作为三乘方的被除数。用三百六十二乘以半通勾的数，得到一千八百五十三万四千四百，作为从方。通勾乘以和步，得到四万八千三百二十，作为第一从廉。通勾的五倍得到一千六百，作为第二从廉。二分五厘作为常法。按带从方廉三乘方法开方，得到八十作为小差。小差就是通股弦差。用通勾减去它，就是城圆的直径。

\switchcolumn[0]*
带从方廉负隅，单位开三乘方曰：置所得三乘方实，以廉隅约之。商得八十，置一于左上为法，置一乘从一廉得三百八十六万五千六百，置一自之以乘从二廉得一千零二十四万，置一自乘再得五十一万二千，以二分五厘因之，得十二万八千为隅法。并从方、一廉、二廉、隅法，得三千二百七十六万八千为下法，与上法相乘，除实尽。

\switchcolumn[1]
\noindent\textbf{【带从方廉负隅开三乘方详解】}\quad 把求得的三乘方被除数，用从廉和隅来约。商得八十，在左上方放一个一作为上法，用一乘以第一从廉得到三百八十六万五千六百，一自乘再乘以第二从廉得到一千零二十四万，一自乘三次得到五十一万二千，用二分五厘乘之得到十二万八千作为隅法。合并从方、第一从廉、第二从廉、隅法，共三千二百七十六万八千作为下法，与上法相乘，把被除数刚好除尽。

\vspace{0.5em}

\Question{问四}{%
东门外往南有树，乙出东门往东不知步数而立。甲出北门东行二百步，斜望乙与树，正与城相参直。既而乙复折而斜行至树下，与甲相望，计乙直行斜行共五十步。
}

\switchcolumn[0]*
\noindent\textbf{释曰}\quad 此以底勾与边勾弦和，立法测望。甲出北门东行，底勾也；乙一直一斜，边勾、边弦也。

\switchcolumn[1]
\noindent\textbf{【释】}\quad 这是用底勾与边勾弦和来建立方法进行测望的问题。甲出北门向东行走，是底勾；乙先直走再斜走，是边勾与边弦。

\switchcolumn[0]*
\noindent\textbf{术曰}\quad 底勾与和相减，余一百五十为差。差加底勾，复以差乘之，得数半之，得二万六千二百五十。差自之，得二万二千五百。二数相减，余三千七百五十为实。并勾和，半之，得一百二十五为法。实如法而得一。

\switchcolumn[1]
\noindent\textbf{【术】}\quad 底勾与和相减，剩余一百五十作为差。差加上底勾，再用差来乘，得到的数取一半，得到二万六千二百五十。差自乘，得到二万二千五百。两个数相减，剩余三千七百五十作为被除数。勾与和相加取一半，得到一百二十五作为除数。被除数除以除数得到一个数值。

\vspace{0.5em}

\Question{问五}{%
南门外往东不知步数有树，乙出南门南行不知步数而立。甲出北门东行二百步，见树与乙与城相参直，乙复斜行至树下与甲相望，计乙一直一斜共二百八十八步。问城径。
}

\switchcolumn[0]*
\noindent\textbf{释曰}\quad 此以底勾与明股弦和，立法测望。甲出北门东行，底勾也；乙出南门南行，明股也；斜行，明弦也。

\switchcolumn[1]
\noindent\textbf{【释】}\quad 这是用底勾与明股弦和来建立方法进行测望的问题。甲出北门向东行走，是底勾；乙出南门向南行走，是明股；乙斜着行走，是明弦。

\switchcolumn[0]*
\noindent\textbf{术曰}\quad 勾和相减余，半之得四十四为半差。以减底勾，余一百五十六为泛率。泛率自之，又倍之，得四万八千六百七十二。半差乘和步，得一万二千六百七十二。二数相减，余三万六千为实。半底勾减和步，得一百八十八。倍泛率，得三百一十二。二数相并，得五百为法。实如法而一，得明勾。

\switchcolumn[1]
\noindent\textbf{【术】}\quad 勾与和相减的余数，取一半得到四十四作为半差。用底勾减去半差，余下一百五十六作为泛率。泛率自乘，再乘以二，得到四万八千六百七十二。半差乘以和步，得到一万二千六百七十二。两个数相减，余下三万六千作为被除数。底勾的一半减去和步，得到一百八十八。泛率的二倍得到三百一十二。两个数相加，得到五百作为除数。被除数除以除数，得到明勾。

%----------------------------------------------------------------------
\Section{勾与较测望二}
%----------------------------------------------------------------------

\Question{问六}{%
甲乙俱在城外西北乾隅，甲南行不知步数而立，乙东行三百二十步见之，甲又斜行与乙相会，计甲直行不及斜行八十步。
}

\switchcolumn[0]*
\noindent\textbf{释曰}\quad 此以通勾与股弦较测望。乙东行，通勾也；甲直行不及斜行，股弦较也。

\switchcolumn[1]
\noindent\textbf{【释】}\quad 这是用通勾与股弦差来进行测望的问题。乙向东行走，是通勾；甲直行比斜行少走的步数，就是股弦差。

\switchcolumn[0]*
\noindent\textbf{术曰}\quad 较除勾算，得一千二百八十为股弦和。减较，半之为股；加较，半之为弦。

\switchcolumn[1]
\noindent\textbf{【术】}\quad 用差来除勾的数，得到一千二百八十作为股弦和。减去差取一半就是股；加上差取一半就是弦。

\switchcolumn[0]*
边勾以下，俱即此类推。

\switchcolumn[1]
边勾以下的各类问题，都按照这个方法类推。

%----------------------------------------------------------------------
\Section{股与和测望三}
%----------------------------------------------------------------------

\Question{问七}{%
甲乙二人俱在城外西北乾隅，甲南行六百步而立，乙东行不知步数见之，又斜行与甲相会，计乙直斜共行一千步。问城径。
}

\switchcolumn[0]*
\noindent\textbf{释曰}\quad 此以通股、勾弦和测望。甲南行，通股也；乙直东行与斜行共，勾弦和也。

\switchcolumn[1]
\noindent\textbf{【释】}\quad 这是用通股与勾弦和来进行测望的问题。甲向南行走，是通股；乙直向东行与斜行共走的步数，是勾弦和。

\switchcolumn[0]*
\noindent\textbf{术曰}\quad 股自之，得三十六万。和除之，得三百六十为勾弦较。减和，半之为勾；加和，半之为弦。

\switchcolumn[1]
\noindent\textbf{【术】}\quad 股自乘，得到三十六万。用和来除，得到三百六十作为勾弦差。用和减去差取一半就是勾；用和加上差取一半就是弦。

\switchcolumn[0]*
边股以下，推此。

\switchcolumn[1]
边股以下的各类问题，都按此方法推算。

\vspace{0.5em}

\Question{问八}{%
甲从乾隅南行六百步而立，乙出南门直行，丙出东门直行，三人相望，俱与城相参直。计其行步，则乙与丙共行一百五十一步。
}

\switchcolumn[0]*
\noindent\textbf{释曰}\quad 此以通股、边勾、明股和，立法测望。甲行，通股；乙行，明股；丙行，边勾也。共之和也。

\switchcolumn[1]
\noindent\textbf{【释】}\quad 这是用通股、边勾、明股和来建立方法进行测望的问题。甲行走的是通股；乙行走的是明股；丙行走的是边勾。三人行走数的总和就是和。

\switchcolumn[0]*
\noindent\textbf{术曰}\quad 通股为算，半而自之，得三百二十四亿，为三乘方实。倍和加通股，以乘半通股算，得一亿六千二百三十六万，为从方。通股乘和步，得九万零六百，为从一廉。通股加半股，得九百，为从二廉。二分五厘为隅算。作带从方廉负隅，以二廉减从，翻法开三乘方法除之，得三百六十为股圆差。以减通股，即圆径。

\switchcolumn[1]
\noindent\textbf{【术】}\quad 以通股为基数，取一半后自乘，得到三百二十四亿，作为三乘方的被除数。二倍和加上通股，再乘以半通股的数，得到一亿六千二百三十六万，作为从方。通股乘以和步，得到九万零六百，作为第一从廉。通股加上半股，得到九百，作为第二从廉。二分五厘作为隅算。按带从方廉负隅的方法，用第二从廉减从方，翻法开三乘方来除之，得到三百六十作为股圆差。用通股减去它，就是圆的直径。

\switchcolumn[0]*
带从方廉负隅，以二廉减从，翻法开三乘方曰：置所得三乘方实，以从方廉隅约之。初商三百，置一于左上为法，置一自之以乘二廉，得八千一百万，以减从方，余八千一百三十六万。置一乘从一廉，得二千七百一十八万。置一自乘再乘，以隅算二分五厘因之，得六百七十五万为隅法。并从方、从一廉、隅法，共一亿一千五百二十九万为下法。与上法相乘，除实三百四十五亿八千七百万，实不满法，反减实三百二十四亿，余二十一亿八千七百万为负积。四因隅法，得二千七百万为方法。初商自之，六因又以隅因之，得一十三万五千为上廉。初商四之，隅因之，得三百为下廉。商次位得六十，置一于左，次为上法，倍初商加次商，得六百六十，以乘从二廉，得五十九万四千。又并初次商，得三百六十，因得二亿四千三百八十四万，以减余从，亦不及减，反减从八千一百三十六万，余一亿三千二百四十八万为负从。置一倍初商加次商，得六百六十，以乘从一廉，得五千九百七十九万六千。置一乘上廉，得八十一万。置一自之以乘下廉，得一百零八万。置一自乘再乘，隅因之，得五万四千为隅法。并方法、从一廉、上下廉、隅法，共九千六百零三万。以减负从，余三千六百四十五万，与上次法，除负积二十一亿八千七百万。

\switchcolumn[1]
\noindent\textbf{【带从方廉负隅翻法开三乘方详解】}\quad 把求得的三乘方被除数，用从方、从廉和隅来约。初次商得三百，在左上方放一个一作为上法，一自乘再乘以第二从廉，得到八千一百万，用来减从方，余下八千一百三十六万。用一乘以第一从廉，得到二千七百一十八万。一自乘三次，用隅算二分五厘乘之，得到六百七十五万作为隅法。合并从方、第一从廉、隅法，共一亿一千五百二十九万作为下法。与上法相乘，从被除数中减去三百四十五亿八千七百万，被除数不够减，反过来减被除数三百二十四亿，余下二十一亿八千七百万作为负积。将隅法四倍，得到二千七百万作为方法。初商自乘，六倍再用隅算乘之，得到十三万五千作为上廉。初商四倍再用隅算乘之，得到三百作为下廉。第二次商得六十，在左上方放一个一作为上法，二倍初商加次商得到六百六十，用来乘第二从廉得到五十九万四千。又合并初次商得三百六十，用它去乘余下从方，也不够减，反过来减从方八千一百三十六万，余下一亿三千二百四十八万作为负从。用一倍初商加次商得到六百六十，去乘第一从廉，得到五千九百七十九万六千。用一去乘上廉，得到八十一万。一自乘再乘下廉，得到一百零八万。一自乘三次再用隅算乘之，得到五万四千作为隅法。合并方法、第一从廉、上廉下廉、隅法，共九千六百零三万。用来减负从，余下三千六百四十五万，与上法相乘，用来除负积二十一亿八千七百万。

%----------------------------------------------------------------------
\Section{弦与和测望四}
%----------------------------------------------------------------------

\Question{问九}{%
甲丙二人俱在城外东北艮隅，乙南行过城门而立，甲东行望乙与城相参直而止。丙丁二人俱在城外西南坤隅，丁向东过城门而立，丙向南行望丁及甲乙，悉与城相参直，丙复斜行六百八十步与甲相会。计乙之南与丁之东共三百四十二步。问城径。
}

\switchcolumn[0]*
\noindent\textbf{释曰}\quad 此通弦与大差勾、小差股和，立法测望。乙从艮隅而南过城门而立，山之艮，小差股也；以甲东行为勾，丁从坤隅东行过城门而立，坤之月，大差勾也；以丙南行为股，丙斜行与甲相会，通弦也。乙丁直行共步，大差勾与小差股和也。

\switchcolumn[1]
\noindent\textbf{【释】}\quad 这是用通弦与大差勾、小差股和来建立方法进行测望的问题。乙从艮隅向南走过城门后停下，艮对应山，是小差股；甲向东行走作为勾，丁从坤隅向东行走走过城门后停下，坤对应月，是大差勾；丙向南行走作为股，丙斜着行走与甲相会，是通弦。乙向南和丁向东的直行共步，就是大差勾与小差股的和。

\switchcolumn[0]*
\noindent\textbf{术曰}\quad 斜步、共步相乘，倍之，得四十六万五千一百二十为实。斜步、共步相减，余三百三十八为差。倍斜行加差，共一千六百九十八为从。作带从开平方法除之，得全径。

\switchcolumn[1]
\noindent\textbf{【术】}\quad 斜行步数与共步相乘，再乘以二，得到四十六万五千一百二十作为被除数。斜行步数与共步相减，余下三百三十八作为差。二倍斜行加上差，共一千六百九十八作为从方。按带从开平方法除之，得到圆的全径。

\switchcolumn[0]*
带从开平方法见前。

\switchcolumn[1]
带从开平方法见前面各卷的讲解。

\vspace{0.5em}

\Question{问十}{%
甲出东门东行，乙出南门南行，各不知步数，相望与城相参直，甲复斜行二百八十九步与乙相会。乙直长，甲直短，共计一百五十一步。问城径。
}

\switchcolumn[0]*
\noindent\textbf{释曰}\quad 此以皇极弦、边勾、明股和，立法测望。甲东行为边勾，乙南行为明股，甲之斜行，皇极弦也。

\switchcolumn[1]
\noindent\textbf{【释】}\quad 这是用皇极弦、边勾、明股和来建立方法进行测望的问题。甲向东行走是边勾，乙向南行走是明股，甲的斜行是皇极弦。

\switchcolumn[0]*
\noindent\textbf{术曰}\quad 斜行自之，得八万三千五百二十一为弦算。共步自之，得二万二千八百零一为和算。和算减弦算，余六万零七百二十为实。倍共步减斜行，余十三步为从。作带从开平方法除之，得全径。

\switchcolumn[1]
\noindent\textbf{【术】}\quad 斜行自乘，得到八万三千五百二十一作为弦数。共步自乘，得到二万二千八百零一作为和数。用和数减弦数，余下六万零七百二十作为被除数。二倍共步减去斜行，余下十三步作为从方。按带从开平方法除之，得到圆的全径。

\switchcolumn[0]*
带从开平方法见前。

\switchcolumn[1]
带从开平方法见前面各卷的讲解。

\vspace{0.5em}

\Question{问十一}{%
甲乙二人同出东门，甲东行乙南行；丙丁二人同出南门，丙南行丁东行，各不知步数而立。四人遥相望，悉与城相参直。问其步数，则曰：甲丙共行了一百五十一步，乙丁立处相距一百零二步。问城径。
}

\switchcolumn[0]*
\noindent\textbf{释曰}\quad 此太虚弦与边勾、明股和，立法测望。甲出东门直行为边勾而乙南行为股，丙出南门南行为明股而丁东行为勾。甲丙共步，边勾、明股和也。乙丁相距，太虚弦也。

\switchcolumn[1]
\noindent\textbf{【释】}\quad 这是用太虚弦与边勾、明股和来建立方法进行测望的问题。甲出东门直走是边勾而乙向南走是股，丙出南门向南走是明股而丁向东走是勾。甲与丙共行的步数，就是边勾与明股的和。乙丁立处之间的距离，就是太虚弦。

\switchcolumn[0]*
\noindent\textbf{术曰}\quad 共步、相距步相减，余四十九为差，自之得二千四百零一为差算。共步自之，得二万二千八百零一为和算。差算减和算，余二万零四百为实。倍距步减差，余一百五十五为从。作以从减法，开平方法除之，得全径。

\switchcolumn[1]
\noindent\textbf{【术】}\quad 共步与相距步相减，余下四十九作为差，差自乘得到二千四百零一作为差数。共步自乘，得到二万二千八百零一作为和数。用差数减和数，余下二万零四百作为被除数。二倍距步减去差，余下一百五十五作为从方。按以从减法的开平方法除之，得到圆的全径。

\switchcolumn[0]*
以从减法开平方法见前。

\switchcolumn[1]
以从减法开平方法见前面各卷的讲解。

\switchcolumn[0]*
又为以从添积开平方。

\switchcolumn[1]
又可以用以从添积的开平方法来求解。

\switchcolumn[0]*
其法曰：初商二百，置一于左上为法，置一乘从，得三万二千为益积，添入原积，共五万一千四百为实。置一为隅法，与上法相乘，除实四万，余实一万一千四百。倍隅法，得四百为廉法。次商四十，置一于左上为法，置一乘从方，得六千二百为益实，添入余积，共一万七千六百为实。置一并廉法，共四百四十为下法，与上法相乘，除实尽。

\switchcolumn[1]
\noindent\textbf{【以从添积开平方详解】}\quad 初次商得二百，在左上方放一个一作为上法，用一乘以从方得到三万二千作为益积，添入原来的积数，共五万一千四百作为被除数。放一个一作为隅法，与上法相乘，从被除数中减去四万，余下一万一千四百。将隅法加倍，得到四百作为廉法。第二次商得四十，在左上方放一个一作为上法，用一乘以从方得到六千二百作为益实，添入剩余的积数，共一万七千六百作为被除数。将一与廉法合并，共四百四十作为下法，与上法相乘，把被除数刚好除尽。

\switchcolumn[0]*
后凡言以从添积开平方法，俱仿此。

\switchcolumn[1]
以后凡提到以从添积开平方法的，都仿照这个方法。

\vspace{0.5em}

\Question{问十二}{%
出南门向东有槐树，出东门向南有柳树。丙丁俱出南门，丙直往南，丁往东至槐树下立；甲乙俱出东门，甲直往东，乙往南至柳树下立。四人遥相望见，各不知步数。只云：丙丁共行了二百零七步，甲乙共行了四十六步，其甲丙立处相距二百八十九步。问城径。
}

\switchcolumn[0]*
\noindent\textbf{释曰}\quad 此以皇极弦与明勾股和、边勾股和，立法测望。槐在南门之东，为南之月，明勾也；丁直行往南，为日之南，明股也；共行二百零七，明勾股和也。柳在东门之南，为山之东，边股也；甲直行往东，为东之川，边勾也；共行四十六步，边勾股和也。甲丙立处相距，为日川，皇极弦也。

\switchcolumn[1]
\noindent\textbf{【释】}\quad 这是用皇极弦与明勾股和、边勾股和来建立方法进行测望的问题。槐树在南门东边，对应南之月，是明勾；丁直往南走，对应日之南，是明股；丙丁共走二百零七步，就是明勾股和。柳树在东门南边，对应山之东，是边股；甲直往东走，对应东之川，是边勾；甲乙共走四十六步，就是边勾股和。甲丙立处之间的距离，对应日川，是皇极弦。

\switchcolumn[0]*
\noindent\textbf{术曰}\quad 二和相减余，以减相距余，半之得六十四为平勾。以加二和相减，为平股。相乘为实，平方开之，即半径。

\switchcolumn[1]
\noindent\textbf{【术】}\quad 两个和相减的余数，用相距减去这个余数，取一半得到六十四作为平勾。用这个数加上二和相减的余数，作为平股。两者相乘作为被除数，开平方，就是圆的半径。

\switchcolumn[0]*
\noindent\textbf{又曰}\quad 二和相并，以减相距余，半之得一十八为泛率。加明和为长，加边和为广。长广相乘，得半径算。

\switchcolumn[1]
\noindent\textbf{【又】}\quad 两个和合并，用相距减去合并后的数，取一半得到一十八作为泛率。泛率加上明和作为长，加上边和作为广。长和广相乘，得到半径的数。

\vspace{0.5em}

\Question{问十三}{%
南门之东有槐，东门之南有柳。丙出南门直行，丁出南门东至槐下；甲出东门直行，乙出东门南至柳下。相望俱与城相参直。计丙南丁东共行二百零七步，甲东乙南共行四十六步，其二树相距一百零二步。问城径。
}

\switchcolumn[0]*
\noindent\textbf{释曰}\quad 此与前问同，前以远相距言，此以近相距言。近相距，太虚弦也。以太虚弦与明、叀二和，立法测望。

\switchcolumn[1]
\noindent\textbf{【释】}\quad 这个问题与前一问相同，前一问说的是远距离相距，这一问说的是近距离相距。近距离相距的步数，就是太虚弦。用太虚弦与明勾股和、叀勾股和来建立方法进行测望。

\switchcolumn[0]*
\noindent\textbf{术曰}\quad 叀和乘虚弦，又自之，得二千二百零一万四千八百六十四为平实。并二和自之，得六万四千零九为二和算。边和自之，得二千一百一十六为边和算。明和自之，得四万二千八百四十九为明和算。并明和算、叀和算，以减二和算，余一万九千零四十四为益隅。作负隅开平方法除之，得叀弦。倍弦算与和算相减，开其余，得叀勾股较。加和，半之为股；减和，半之为勾。

\switchcolumn[1]
\noindent\textbf{【术】}\quad 用叀和乘以虚弦，再自乘，得到二千二百零一万四千八百六十四作为平实。将二和合并后自乘，得到六万四千零九作为二和数。边和自乘，得到二千一百一十六作为边和数。明和自乘，得到四万二千八百四十九作为明和数。将明和数与叀和数合并，用二和数来减，余下一万九千零四十四作为益隅。按负隅开平方法除之，得到叀弦。二倍弦数与和数相减，开其余数，得到叀勾股差。加上和取一半是股；减去和取一半是勾。

\switchcolumn[0]*
负隅开平方曰：置所得平实，以益隅约之。初商三十，置一于左上为法，置一乘益隅得五十七万一千三百二十为下法，与上法相乘，除实一千七百一十三万九千六百，余实四百八十七万五千二百六十四。倍下法，得一百一十四万二千六百四十为廉法。约次商得四，置一于左上为法，置一乘益隅得七万六千一百七十六，并入廉法，共一百二十一万八千八百一十六为下法，与上法相乘，除实尽。

\switchcolumn[1]
\noindent\textbf{【负隅开平方详解】}\quad 把求得的平实，用益隅来约。初次商得三十，在左上方放一个一作为上法，用一乘以益隅得到五十七万一千三百二十作为下法，与上法相乘，从被除数中减去一千七百一十三万九千六百，余下四百八十七万五千二百六十四。将下法加倍，得到一百一十四万二千六百四十作为廉法。约得第二次商为四，在左上方放一个一作为上法，用一乘以益隅得到七万六千一百七十六，并入廉法，共一百二十一万八千八百一十六作为下法，与上法相乘，把剩余被除数刚好除尽。

\switchcolumn[0]*
此法已见一卷底勾弦条下，因隅算多故重出。

\switchcolumn[1]
这个方法在第一卷底勾弦条下已经出现过，因为这里的隅算较多所以重新列出。

\switchcolumn[0]*
\noindent\textbf{又曰}\quad 隅算除平实，即得叀弦算。

\switchcolumn[1]
\noindent\textbf{【又】}\quad 用平实除以隅算，就直接得到叀弦的数值。

\switchcolumn[0]*
\noindent\textbf{又曰}\quad 明和乘虚弦，又自之，得四亿四千五百八十万零九百九十六为平实。如前法为负隅，平方开之，得明弦。若以益隅除平实，径得明弦算。

\switchcolumn[1]
\noindent\textbf{【又】}\quad 用明和乘以虚弦，再自乘，得到四亿四千五百八十万零九百九十六作为平实。按照前面的方法设负隅，开平方，得到明弦。如果用益隅来除平实，就直接得到明弦的数值。

\switchcolumn[0]*
\noindent\textbf{又术}\quad 虚弦自之，得一万零四百零四为虚弦算。以叀和乘之，得四十七万八千五百八十四为平实。倍明和，得四百一十四为益隅，开之得叀弦。若以益隅除平实，径得叀弦算。

\switchcolumn[1]
\noindent\textbf{【又术】}\quad 虚弦自乘，得到一万零四百零四作为虚弦数。用叀和乘之，得到四十七万八千五百八十四作为平实。二倍明和得到四百一十四作为益隅，开平方得到叀弦。如果用益隅来除平实，就直接得到叀弦的数值。

\switchcolumn[0]*
虚弦自之，以明和乘之，得二百一十五万三千六百二十八为平实。倍叀和为益隅，开之得明弦。若以益隅除平实，径得明弦算。

\switchcolumn[1]
虚弦自乘，用明和乘之，得到二百一十五万三千六百二十八作为平实。二倍叀和作为益隅，开平方得到明弦。如果用益隅来除平实，就直接得到明弦的数值。

\switchcolumn[0]*
三位负隅开平方曰：置平实四亿四千五百八十万零九百九十六于左，以益隅一万九千零四十四约之。初商一百，置一于左上为法，置一于右下乘益隅，得一百九十万四千四百为下法，与上法相乘，除实一亿九千零四十四万，余实二亿五千五百三十六万零九百九十六。倍下法，得三百八十万八千八百为廉法。次商五十，置一于左上为法，置一乘益隅，得九十五万二千二百为隅法，并廉法，共四百七十六万一千为下法，与上次相乘，除实二亿三千八百零五万，余实一千七百三十一万零九百九十六。倍隅法，得一百九十万四千四百，并入廉法，共五百七十一万三千二百为廉法。约三商得三，置一于左为法，置一右下乘益隅，得五万七千一百三十二为隅法，并入廉法，共五百七十七万零三百三十二为下法，与上法相乘，除实尽。

\switchcolumn[1]
\noindent\textbf{【三位负隅开平方详解】}\quad 把平实四亿四千五百八十万零九百九十六放在左边，用益隅一万九千零四十四来约。初次商得一百，在左上方放一个一作为上法，在右下方放一个一乘益隅，得到一百九十万四千四百作为下法，与上法相乘，从被除数中减去一亿九千零四十四万，余下二亿五千五百三十六万零九百九十六。将下法加倍，得到三百八十万八千八百作为廉法。第二次商得五十，在左上方放一个一作为上法，用一乘益隅得到九十五万二千二百作为隅法，并入廉法，共四百七十六万一千作为下法，与上法相乘，减去二亿三千八百零五万，余下一千七百三十一万零九百九十六。将隅法加倍，得到一百九十万四千四百，并入廉法，共五百七十一万三千二百作为廉法。约得第三次商为三，在左方放一个一作为上法，在右下放一个一乘益隅得到五万七千一百三十二作为隅法，并入廉法，共五百七十七万零三百三十二作为下法，与上法相乘，把被除数刚好除尽。

%----------------------------------------------------------------------
\Section{弦与较测望六}
%----------------------------------------------------------------------

\Question{问十四}{%
甲丙二人俱在城外西北隅起程，丙南行甲东行，各不知步数，隔城相望。既而甲斜行六百八十步与丙相会。问其东行步数，则曰：我少于丙南行二百八十步。问城径。
}

\switchcolumn[0]*
\noindent\textbf{释曰}\quad 此通弦与通勾股较，立法测望。甲东行为勾，丙南行为股，甲少于丙步数，勾股较也，斜行，弦也。

\switchcolumn[1]
\noindent\textbf{【释】}\quad 这是用通弦与通勾股差来建立方法进行测望的问题。甲向东行走作为勾，丙向南行走作为股，甲比丙少走的步数，就是勾股差，斜走的步数是弦。

\switchcolumn[0]*
\noindent\textbf{术曰}\quad 弦自乘，倍之，得九十二万四千八百。较自乘，得七万八千四百。相减，余八十四万六千四百为实。平方开之，得勾股和九百二十。加较，半之为股；减较，半之为勾。

\switchcolumn[1]
\noindent\textbf{【术】}\quad 弦自乘，再乘以二，得到九十二万四千八百。差自乘，得到七万八千四百。两数相减，余下八十四万六千四百作为被除数。开平方，得到勾股和九百二十。加上差取一半是股；减去差取一半是勾。

\switchcolumn[0]*
\noindent\textbf{又曰}\quad 弦较相减，得四百为弦较较；相并，得九百六十为弦较和。弦较较、弦较和相乘，得三十八万四千为实。倍较，得五百六十为从，二为隅算。作以从减法、负隅开平方法除之，得通股。作带从负隅开平方法除之，得通勾。

\switchcolumn[1]
\noindent\textbf{【又】}\quad 弦与差相减，得到四百作为弦较较；两者相加，得到九百六十作为弦较和。弦较较与弦较和相乘，得到三十八万四千作为被除数。二倍差得到五百六十作为从方，二作为隅算。按以从减法、负隅开平方法除之，得到通股。按带从负隅开平方法除之，得到通勾。

\switchcolumn[0]*
带从负隅开平方法，见四卷底勾通弦条。

\switchcolumn[1]
带从负隅开平方法，见第四卷底勾通弦条。

\switchcolumn[0]*
带从负隅以从减隅开平方法，见四卷大差勾黄长弦条下。

\switchcolumn[1]
带从负隅以从减隅开平方法，见第四卷大差勾黄长弦条下。

\switchcolumn[0]*
又为以从添积负隅开平方，以六百乘从益实，倍六百，得一千二百为法。即是边弦以下类推。

\switchcolumn[1]
又可以用以从添积负隅开平方的方法，用六百乘从方作为益实，二倍六百得到一千二百作为法。边弦以下各类问题都可以类推。

\vspace{0.5em}

\Question{问十五}{%
乙出东门南行不知步数而立，甲出西门直往南行，回望乙与城相参直，又斜行五百一十步与乙相会。问乙行步，则曰：少于城径二百一十步。不知城径几何。
}

\switchcolumn[0]*
\noindent\textbf{释曰}\quad 此黄广弦与叀股、黄广勾较，立法测望。乙出东门南行为叀股，城径即黄广勾，少于城径即叀股、黄广勾较也，斜行，黄广弦也。

\switchcolumn[1]
\noindent\textbf{【释】}\quad 这是用黄广弦与叀股、黄广勾差来建立方法进行测望的问题。乙出东门向南行走是叀股，城圆的直径就是黄广勾，比城径少的步数就是叀股与黄广勾的差，斜走的步数是黄广弦。

\switchcolumn[0]*
\noindent\textbf{术曰}\quad 较自之，得四万四千一百为较算，以为实。斜步四之，减二较，余一千六百二十为从，五为隅算。作负隅减从，开平方法除之，得叀股三十，加较为黄广勾，即城径。

\switchcolumn[1]
\noindent\textbf{【术】}\quad 差自乘，得到四万四千一百作为差数，作为被除数。斜步的四倍减去二倍差，余下一千六百二十作为从方，五作为隅算。按负隅减从的开平方法除之，得到叀股三十，加上差就是黄广勾，也就是城圆的直径。

\switchcolumn[0]*
负隅减从开平方法，见二卷通勾叀勾条。

\switchcolumn[1]
负隅减从开平方法，见第二卷通勾叀勾条。

\vspace{0.5em}

\Question{问十六}{%
乙出南门东行不知步数而立，甲出北门直往东行，望乙与城相参直，又斜行二百七十二步与乙相会。问乙东行步，则曰：少于城径一百六十八步。不知城径几何。
}

\switchcolumn[0]*
\noindent\textbf{释曰}\quad 此黄长弦与明勾、黄长股较，立法测望。乙出南门东行为明勾，城径即黄长股，少于城径即明勾、黄长股较也，斜行，黄长弦也。

\switchcolumn[1]
\noindent\textbf{【释】}\quad 这是用黄长弦与明勾、黄长股差来建立方法进行测望的问题。乙出南门向东行走是明勾，城圆的直径就是黄长股，比城径少的步数就是明勾与黄长股的差，斜走的步数是黄长弦。

\switchcolumn[0]*
\noindent\textbf{术曰}\quad 较自之，得二万八千二百二十四为实。四斜行减二较，余七百五十二为从方，五为隅算。作负隅减从，开平方法除之，得明勾七十二，加较为黄长股，即城径。

\switchcolumn[1]
\noindent\textbf{【术】}\quad 差自乘，得到二万八千二百二十四作为被除数。四倍斜行减去二倍差，余下七百五十二作为从方，五作为隅算。按负隅减从的开平方法除之，得到明勾七十二，加上差就是黄长股，也就是城圆的直径。

\switchcolumn[0]*
负隅减从开平方法，见二卷。

\switchcolumn[1]
负隅减从开平方法，见第二卷。

\EndVol{测圆海镜分类释术卷六终}

%======================================================================
\Volume{七}
%======================================================================

\switchcolumn[0]*
\noindent\textbf{测圆海镜分类释术卷七}\quad 元~李~冶~撰\\
明~顾应祥~释术

\switchcolumn[1]
\noindent\textbf{测圆海镜分类释术·第七卷}\quad 元代~李~冶~撰\\
明代~顾应祥~释术

\vspace{1em}

\switchcolumn[0]*
卷七专论大差、小差、中差之测望。大差者，通弦与通股之差；小差者，通弦与通勾之差；中差者，大小差之差也。此三差为测圆术之枢纽，诸题皆以之为本。凡九题，分为三节。

\switchcolumn[1]
第七卷专门论述大差、小差、中差的测望方法。大差，就是通弦与通股的差；小差，就是通弦与通勾的差；中差，就是大差与小差的差。这三种差是测圆术的关键，所有题目都以它们为基础。共有九道题，分为三节。

%----------------------------------------------------------------------
\Section{大差测望一}
%----------------------------------------------------------------------

\Question{问一}{%
甲乙二人俱在城外西北乾隅，甲南行六百步而立，乙东行三百二十步而立。既而甲斜行与乙相会，计甲共行一千二百八十步。以甲斜行与乙东行相较，求大差。
}

\switchcolumn[0]*
\noindent\textbf{释曰}\quad 此以大差立法测望。甲南行六百步为通股，乙东行三百二十步为通勾。甲斜行为通弦。通弦减通股，余为大差。大差勾者，以通勾与大差相减之数也。

\switchcolumn[1]
\noindent\textbf{【释】}\quad 这是用大差来建立方法进行测望的问题。甲向南行走六百步是通股，乙向东行走三百二十步是通勾。甲斜着行走是通弦。通弦减去通股，余下就是大差。大差勾，就是用通勾减去大差后得到的数。

\switchcolumn[0]*
\noindent\textbf{术曰}\quad 通弦自乘，以通股自乘减之，开方得大差勾。以勾股求之，得城径。

\switchcolumn[1]
\noindent\textbf{【术】}\quad 通弦自乘，减去通股自乘，开平方得到大差勾。再用勾股术推算，求得城圆的直径。

\vspace{0.5em}

\Question{问二}{%
甲从乾隅东行三百二十步而立，乙从坤隅南行六百步而立。二人隔城相望，甲复斜行至乙处，计甲共行一千二百八十步。以乙南行与甲斜行相较，求大差。
}

\switchcolumn[0]*
\noindent\textbf{释曰}\quad 此亦以大差立法测望。通股减通弦，余为大差。大差勾即通勾之余数。

\switchcolumn[1]
\noindent\textbf{【释】}\quad 这也是用大差来建立方法进行测望的问题。通股减去通弦，余下就是大差。大差勾就是通勾的余数。

\switchcolumn[0]*
\noindent\textbf{术曰}\quad 通弦自乘，减通股自乘，开方得大差。以大差减通股，得股圆差，即城径。

\switchcolumn[1]
\noindent\textbf{【术】}\quad 通弦自乘，减去通股自乘，开平方得到大差。用大差减通股，得到股圆差，就是城圆的直径。

\vspace{0.5em}

\Question{问三}{%
乙出东门南行，不知步数而立。甲从乾隅东行三百二十步，望乙与城相参直。甲又斜行五百一十步与乙相会。求大差勾、股、弦。
}

\switchcolumn[0]*
\noindent\textbf{释曰}\quad 此以明勾与大差股测望。甲东行，通勾也；乙出东门南行，大差股也；甲斜行，大差弦也。

\switchcolumn[1]
\noindent\textbf{【释】}\quad 这是用明勾与大差股来进行测望的问题。甲向东行走，是通勾；乙出东门向南行走，是大差股；甲斜着行走，是大差弦。

\switchcolumn[0]*
\noindent\textbf{术曰}\quad 通勾自乘，以明勾减之，余为实。以大差弦除之，得大差股。以大差股减通股，余即城径。

\switchcolumn[1]
\noindent\textbf{【术】}\quad 通勾自乘，减去明勾，余下作为被除数。用大差弦来除，得到大差股。用大差股减通股，余下就是城圆的直径。

\vspace{0.5em}

\Question{问四}{%
乙出南门东行，不知步数而立。甲从乾隅南行六百步，望乙与城相参直。甲又斜行至乙处，计甲共行六百八十步。求大差。
}

\switchcolumn[0]*
\noindent\textbf{释曰}\quad 此以大差股与通弦测望。大差股即通股减弦之余。

\switchcolumn[1]
\noindent\textbf{【释】}\quad 这是用大差股与通弦来进行测望的问题。大差股就是通股减去弦后的余数。

\switchcolumn[0]*
\noindent\textbf{术曰}\quad 通弦自乘，减通股自乘，开方得大差勾。以大差勾减通勾，余为城径。

\switchcolumn[1]
\noindent\textbf{【术】}\quad 通弦自乘，减去通股自乘，开平方得到大差勾。用大差勾减通勾，余下就是城圆的直径。

\vspace{0.5em}

\Question{问五}{%
乙出东门东行，不知步数而立。丙出南门南行，不知步数而立。甲从乾隅东行三百二十步，望乙丙与城相参直，计乙丙共行一百五十一步。以乙行与丙行相较，求大差。
}

\switchcolumn[0]*
\noindent\textbf{释曰}\quad 此以明勾、大差股和测望。乙东行为边勾，丙南行为大差股，共行一百五十一步，即大差勾股和也。

\switchcolumn[1]
\noindent\textbf{【释】}\quad 这是用明勾、大差股和来进行测望的问题。乙向东行走是边勾，丙向南行走是大差股，共走一百五十一步，就是大差勾股和。

\switchcolumn[0]*
\noindent\textbf{术曰}\quad 共步自乘，以通勾乘之为实。以通股减共步为从，开平方得大差股。以减通股即城径。

\switchcolumn[1]
\noindent\textbf{【术】}\quad 共步自乘，用通勾来乘作为被除数。用通股减共步作为从方，开平方得到大差股。用通股减去它就是城圆的直径。

\vspace{0.5em}

\Question{问六}{%
甲从坤隅东行过城门而止，乙从艮隅南行过城门而止。二人相望与城相参直，计甲东行与乙南行共三百四十二步。以甲行减乙行，余八十步。求大差。
}

\switchcolumn[0]*
\noindent\textbf{释曰}\quad 此以大差勾股和与较测望。甲东行过城门为大差勾，乙南行过城门为大差股。

\switchcolumn[1]
\noindent\textbf{【释】}\quad 这是用大差勾股和与差来进行测望的问题。甲向东行走走过城门是大差勾，乙向南行走走过城门是大差股。

\switchcolumn[0]*
\noindent\textbf{术曰}\quad 和自乘，减较自乘，以四因之为实。倍和为从，开平方得大差弦。以弦减通弦，余即城径。

\switchcolumn[1]
\noindent\textbf{【术】}\quad 和自乘，减去差自乘，再乘以四作为被除数。二倍和作为从方，开平方得到大差弦。用通弦减去它，余下就是城圆的直径。

\vspace{0.5em}

\Question{问七}{%
甲乙二人俱出东门，甲东行乙南行。丙丁二人俱出南门，丙南行丁东行。四人相望，各与城相参直。计甲丙共行二百步，乙丁相距二百八十九步。求大差。
}

\switchcolumn[0]*
\noindent\textbf{释曰}\quad 此以太虚弦与边勾明股和测望。乙丁相距即太虚弦，大差之属也。

\switchcolumn[1]
\noindent\textbf{【释】}\quad 这是用太虚弦与边勾明股和来进行测望的问题。乙丁相距的距离就是太虚弦，属于大差这一类。

\switchcolumn[0]*
\noindent\textbf{术曰}\quad 以太虚弦减皇极弦，余为大差弦。以勾股术推之，得城径。

\switchcolumn[1]
\noindent\textbf{【术】}\quad 用太虚弦减皇极弦，余下就是大差弦。用勾股术推算，得到城圆的直径。

\vspace{0.5em}

\Question{问八}{%
甲从乾隅东行三百二十步，南门外有树，乙出南门南行至树下。甲斜望乙与树及城相参直，计甲斜行二百八十九步。求大差股。
}

\switchcolumn[0]*
\noindent\textbf{释曰}\quad 此以明勾与大差弦测望。甲东行为明勾，斜行为大差弦，乙南行为大差股。

\switchcolumn[1]
\noindent\textbf{【释】}\quad 这是用明勾与大差弦来进行测望的问题。甲向东行走是明勾，斜着行走是大差弦，乙向南行走是大差股。

\switchcolumn[0]*
\noindent\textbf{术曰}\quad 明勾自乘，减大差弦自乘，开方得大差股。以减通股，得城径。

\switchcolumn[1]
\noindent\textbf{【术】}\quad 明勾自乘，减去大差弦自乘，开平方得到大差股。用通股减去它，得到城圆的直径。

\vspace{0.5em}

\Question{问九}{%
乙出东门东行二百步而立，丙出南门南行一百五十一步而立。甲从乾隅东行三百二十步，望乙丙与城相参直。求大差勾股较。
}

\switchcolumn[0]*
\noindent\textbf{释曰}\quad 此以边勾、明股与大差测望。大差勾股较即通股弦较与通勾弦较之差。

\switchcolumn[1]
\noindent\textbf{【释】}\quad 这是用边勾、明股与大差来进行测望的问题。大差勾股差就是通股弦差与通勾弦差的差。

\switchcolumn[0]*
\noindent\textbf{术曰}\quad 边勾自乘，明股自乘，相减为实。以边勾明股相乘为法，除之得大差勾股较。以加减术推之，得城径。

\switchcolumn[1]
\noindent\textbf{【术】}\quad 边勾自乘，明股自乘，相减作为被除数。用边勾与明股相乘作为除数，除之得到大差勾股差。再用加减术推算，得到城圆的直径。

%----------------------------------------------------------------------
\Section{小差测望二}
%----------------------------------------------------------------------

\Question{问十}{%
甲乙二人俱在城外西北乾隅，甲南行六百步而立，乙东行三百二十步而立。甲斜行与乙相会，计甲共行一千二百八十步。以通弦减通勾，求小差。
}

\switchcolumn[0]*
\noindent\textbf{释曰}\quad 此以小差立法测望。通弦减通勾之余，为小差。小差股者，通股之减数也。

\switchcolumn[1]
\noindent\textbf{【释】}\quad 这是用小差来建立方法进行测望的问题。通弦减去通勾的余数，就是小差。小差股，就是通股的减数。

\switchcolumn[0]*
\noindent\textbf{术曰}\quad 通弦自乘，减通勾自乘，开方得小差股。以小差股减通股，得城径。

\switchcolumn[1]
\noindent\textbf{【术】}\quad 通弦自乘，减去通勾自乘，开平方得到小差股。用小差股减通股，得到城圆的直径。

\vspace{0.5em}

\Question{问十一}{%
乙出南门东行，不知步数而立。甲从乾隅南行六百步，望乙与城相参直。甲又斜行二百七十二步与乙相会。求小差。
}

\switchcolumn[0]*
\noindent\textbf{释曰}\quad 此以明勾与小差弦测望。乙出南门东行为明勾，甲斜行为小差弦，甲南行为小差股。

\switchcolumn[1]
\noindent\textbf{【释】}\quad 这是用明勾与小差弦来进行测望的问题。乙出南门向东行走是明勾，甲斜着行走是小差弦，甲向南行走是小差股。

\switchcolumn[0]*
\noindent\textbf{术曰}\quad 明勾自乘，减小差弦自乘，开方得小差股。以小差股减通股，得城径。

\switchcolumn[1]
\noindent\textbf{【术】}\quad 明勾自乘，减去小差弦自乘，开平方得到小差股。用小差股减通股，得到城圆的直径。

\vspace{0.5em}

\Question{问十二}{%
甲从艮隅南行过城门而止，乙从坤隅东行过城门而止。二人相望与城相参直，计甲南行三百六十步，乙东行一百六十步。求小差。
}

\switchcolumn[0]*
\noindent\textbf{释曰}\quad 此以小差勾股直接测望。甲南行过城门为小差股，乙东行过城门为小差勾。

\switchcolumn[1]
\noindent\textbf{【释】}\quad 这是用小差勾股直接进行测望的问题。甲向南行走走过城门是小差股，乙向东行走走过城门是小差勾。

\switchcolumn[0]*
\noindent\textbf{术曰}\quad 小差勾股各自乘，相并开方得小差弦。以通弦减之，余即城径。

\switchcolumn[1]
\noindent\textbf{【术】}\quad 小差勾与小差股各自乘，相加后开平方得到小差弦。用通弦减去它，余下就是城圆的直径。

\vspace{0.5em}

\Question{问十三}{%
甲乙二人同出东门，甲东行二百步，乙南行一百五十一步而立。二人相望与城相参直。求小差勾股较。
}

\switchcolumn[0]*
\noindent\textbf{释曰}\quad 此以小差勾股与通弦测望。甲东行为小差勾，乙南行为小差股。

\switchcolumn[1]
\noindent\textbf{【释】}\quad 这是用小差勾股与通弦来进行测望的问题。甲向东行走是小差勾，乙向南行走是小差股。

\switchcolumn[0]*
\noindent\textbf{术曰}\quad 小差勾股各自乘相并开方得小差弦，以通弦减之得小差。小差减通勾得城径。

\switchcolumn[1]
\noindent\textbf{【术】}\quad 小差勾与小差股各自乘，相加后开平方得到小差弦，用通弦减去它得到小差。小差减通勾得到城圆的直径。

%----------------------------------------------------------------------
\Section{中差测望三}
%----------------------------------------------------------------------

\Question{问十四}{%
甲乙二人俱在城外西北乾隅，甲南行六百步，乙东行三百二十步。各不知余数。以通股减通弦为大差，以通勾减通弦为小差。大小差相减为中差。求中差。
}

\switchcolumn[0]*
\noindent\textbf{释曰}\quad 此以中差立法测望。中差者，大差与小差之差，即通股与通勾之差也。

\switchcolumn[1]
\noindent\textbf{【释】}\quad 这是用中差来建立方法进行测望的问题。中差，就是大差与小差的差，也就是通股与通勾的差。

\switchcolumn[0]*
\noindent\textbf{术曰}\quad 通股自乘，减通勾自乘，开方得勾股和。以和减通弦，余为中差。以中差加减术推之，得城径。

\switchcolumn[1]
\noindent\textbf{【术】}\quad 通股自乘，减去通勾自乘，开平方得到勾股和。用和减通弦，余下就是中差。用中差加减术推算，得到城圆的直径。

\vspace{0.5em}

\Question{问十五}{%
乙出东门南行，不知步数而立。甲从乾隅东行三百二十步，望乙与城相参直。计甲斜行五百一十步。以乙行与城径相较，求中差。
}

\switchcolumn[0]*
\noindent\textbf{释曰}\quad 此以明勾与中差测望。明勾即通勾减边勾之余，中差即大差减小差。

\switchcolumn[1]
\noindent\textbf{【释】}\quad 这是用明勾与中差来进行测望的问题。明勾就是通勾减去边勾的余数，中差就是大差减去小差。

\switchcolumn[0]*
\noindent\textbf{术曰}\quad 明勾自乘为实，以通弦约之，得中差股。以减通股得城径。

\switchcolumn[1]
\noindent\textbf{【术】}\quad 明勾自乘作为被除数，用通弦来约，得到中差股。用通股减去它得到城圆的直径。

\vspace{0.5em}

\Question{问十六}{%
甲从乾隅南行六百步，乙从乾隅东行三百二十步。以甲斜行与乙东行相较，求中差勾股较。
}

\switchcolumn[0]*
\noindent\textbf{释曰}\quad 此以通勾股较测望。中差勾股较即通股减通勾之数。

\switchcolumn[1]
\noindent\textbf{【释】}\quad 这是用通勾股差来进行测望的问题。中差勾股差就是通股减去通勾的数。

\switchcolumn[0]*
\noindent\textbf{术曰}\quad 通股减通勾，得三百八十步为中差。以勾股术推之，得城径。

\switchcolumn[1]
\noindent\textbf{【术】}\quad 通股减去通勾，得到三百八十步作为中差。用勾股术推算，得到城圆的直径。

\vspace{0.5em}

\Question{问十七}{%
乙出南门东行七十二步而立，丙出东门南行三十步而立。甲从乾隅东行三百二十步，望乙丙与城相参直。求中差弦。
}

\switchcolumn[0]*
\noindent\textbf{释曰}\quad 此以明勾、叀股与中差测望。乙东行为明勾，丙南行为叀股。

\switchcolumn[1]
\noindent\textbf{【释】}\quad 这是用明勾、叀股与中差来进行测望的问题。乙向东行走是明勾，丙向南行走是叀股。

\switchcolumn[0]*
\noindent\textbf{术曰}\quad 明勾、叀股各自乘相并开方得中差弦。以通弦减之，余为城径。

\switchcolumn[1]
\noindent\textbf{【术】}\quad 明勾与叀股各自乘，相加后开平方得到中差弦。用通弦减去它，余下就是城圆的直径。

\vspace{0.5em}

\Question{问十八}{%
甲从坤隅东行过城门，乙从艮隅南行过城门。二人相望与城相参直，计甲东行三百二十步，乙南行六百步。以甲行减乙行，求中差。
}

\switchcolumn[0]*
\noindent\textbf{释曰}\quad 此以大差勾股与小差勾股测望中差。中差即大小差之差。

\switchcolumn[1]
\noindent\textbf{【释】}\quad 这是用大差勾股与小差勾股来测望中差的问题。中差就是大小差的差。

\switchcolumn[0]*
\noindent\textbf{术曰}\quad 大差股减小差股，得中差。以减通弦，余为城径。

\switchcolumn[1]
\noindent\textbf{【术】}\quad 大差股减去小差股，得到中差。用通弦减去它，余下就是城圆的直径。

\EndVol{测圆海镜分类释术卷七终}

%======================================================================
\Volume{八}
%======================================================================

\switchcolumn[0]*
\noindent\textbf{测圆海镜分类释术卷八}\quad 元~李~冶~撰\\
明~顾应祥~释术

\switchcolumn[1]
\noindent\textbf{测圆海镜分类释术·第八卷}\quad 元代~李~冶~撰\\
明代~顾应祥~释术

\vspace{1em}

\switchcolumn[0]*
卷八专论皇极、太虚、明、叀四种之测望，以及杂题之类。皇极者，通勾股之中；太虚者，虚勾股之极；明者，明勾股之谓；叀者，叀勾股之称。凡三十四题，分为五节。

\switchcolumn[1]
第八卷专门论述皇极、太虚、明、叀四种测望方法，以及杂题之类。皇极，就是通勾股的中间值；太虚，就是虚勾股的极限值；明，就是明勾股的称谓；叀，就是叀勾股的称呼。共有三十四道题，分为五节。

%----------------------------------------------------------------------
\Section{皇极测望一}
%----------------------------------------------------------------------

\Question{问一}{%
甲乙二人俱在城外西北乾隅，甲南行六百步而立，乙东行三百二十步而立。二人隔城相望，各与城相参直。以勾股之中数，求皇极弦。
}

\switchcolumn[0]*
\noindent\textbf{释曰}\quad 此以皇极弦立法测望。皇极者，勾股之中极也。皇极弦即通弦之半，为勾股容圆之枢纽。

\switchcolumn[1]
\noindent\textbf{【释】}\quad 这是用皇极弦来建立方法进行测望的问题。皇极，就是勾股之中极。皇极弦就是通弦的一半，是勾股容圆的关键。

\switchcolumn[0]*
\noindent\textbf{术曰}\quad 通勾股各自乘相并，开方得通弦。半之为皇极弦。以皇极弦减通弦，余即城径。

\switchcolumn[1]
\noindent\textbf{【术】}\quad 通勾与通股各自乘，相加后开平方得到通弦。取一半作为皇极弦。用通弦减去皇极弦，余下就是城圆的直径。

\vspace{0.5em}

\Question{问二}{%
乙出东门东行，不知步数而立。丙出南门南行，不知步数而立。甲从乾隅东行三百二十步，望乙丙与城相参直，计乙丙共行一百五十一步。以边勾明股求皇极。
}

\switchcolumn[0]*
\noindent\textbf{释曰}\quad 此以边勾明股和测望皇极。乙东行为边勾，丙南行为明股，共行一百五十一步，即皇极勾股和也。

\switchcolumn[1]
\noindent\textbf{【释】}\quad 这是用边勾明股和来测望皇极的问题。乙向东行走是边勾，丙向南行走是明股，共走一百五十一步，就是皇极勾股和。

\switchcolumn[0]*
\noindent\textbf{术曰}\quad 边勾明股各自乘相并开方得皇极弦。以通弦减之余为城径。

\switchcolumn[1]
\noindent\textbf{【术】}\quad 边勾与明股各自乘，相加后开平方得到皇极弦。用通弦减去其余数就是城圆的直径。

\vspace{0.5em}

\Question{问三}{%
甲从乾隅东行三百二十步，乙出东门南行一百五十一步。甲望乙与城相参直，又斜行二百八十九步与乙相会。求皇极勾股较。
}

\switchcolumn[0]*
\noindent\textbf{释曰}\quad 此以皇极弦与边勾明股和测望。甲斜行二百八十九步为皇极弦，乙丙共行一百五十一步为边勾明股和。

\switchcolumn[1]
\noindent\textbf{【释】}\quad 这是用皇极弦与边勾明股和来进行测望的问题。甲斜着行走二百八十九步是皇极弦，乙丙共走一百五十一步是边勾明股和。

\switchcolumn[0]*
\noindent\textbf{术曰}\quad 皇极弦自乘，减边勾明股和自乘，开方得皇极勾股较。以加减术推之，得城径。

\switchcolumn[1]
\noindent\textbf{【术】}\quad 皇极弦自乘，减去边勾明股和自乘，开平方得到皇极勾股差。用加减术推算，得到城圆的直径。

\vspace{0.5em}

\Question{问四}{%
甲从艮隅南行过城门而止，乙从坤隅东行过城门而止。二人相望与城相参直。以甲行乙行求皇极弦。
}

\switchcolumn[0]*
\noindent\textbf{释曰}\quad 此以大差勾小差股测望皇极。甲南行过城门为小差股，乙东行过城门为大差勾。

\switchcolumn[1]
\noindent\textbf{【释】}\quad 这是用大差勾小差股来测望皇极的问题。甲向南行走走过城门是小差股，乙向东行走走过城门是大差勾。

\switchcolumn[0]*
\noindent\textbf{术曰}\quad 大差勾小差股各自乘相并开方得皇极弦。以通弦减之得城径。

\switchcolumn[1]
\noindent\textbf{【术】}\quad 大差勾与小差股各自乘，相加后开平方得到皇极弦。用通弦减去它得到城圆的直径。

\vspace{0.5em}

\Question{问五}{%
甲出东门东行二百步，乙出南门南行一百五十一步。二人相望与城相参直，计二人相距二百八十九步。以相距与各行相较，求皇极。
}

\switchcolumn[0]*
\noindent\textbf{释曰}\quad 此以皇极勾股弦直接测望。甲东行为皇极勾，乙南行为皇极股，相距为皇极弦。

\switchcolumn[1]
\noindent\textbf{【释】}\quad 这是用皇极勾股弦直接进行测望的问题。甲向东行走是皇极勾，乙向南行走是皇极股，两人相距是皇极弦。

\switchcolumn[0]*
\noindent\textbf{术曰}\quad 皇极勾股各自乘相并开方得皇极弦。以勾股术求之，得城径。

\switchcolumn[1]
\noindent\textbf{【术】}\quad 皇极勾与皇极股各自乘，相加后开平方得到皇极弦。用勾股术计算，得到城圆的直径。

%----------------------------------------------------------------------
\Section{太虚测望二}
%----------------------------------------------------------------------

\Question{问六}{%
甲乙二人同出东门，甲东行乙南行。丙丁二人同出南门，丙南行丁东行。四人相望，悉与城相参直。计甲丙共行一百五十一步，乙丁相距一百零二步。以太虚弦测望。
}

\switchcolumn[0]*
\noindent\textbf{释曰}\quad 此以太虚弦立法测望。乙丁相距一百零二步，即太虚弦也。太虚者，虚勾股之极，为最小之勾股形。

\switchcolumn[1]
\noindent\textbf{【释】}\quad 这是用太虚弦来建立方法进行测望的问题。乙丁相距一百零二步，就是太虚弦。太虚，就是虚勾股的极限，是最小的勾股形。

\switchcolumn[0]*
\noindent\textbf{术曰}\quad 太虚弦自乘为实，以通弦约之，得太虚勾。以太虚勾减通勾，得城径。

\switchcolumn[1]
\noindent\textbf{【术】}\quad 太虚弦自乘作为被除数，用通弦来约，得到太虚勾。用太虚勾减通勾，得到城圆的直径。

\vspace{0.5em}

\Question{问七}{%
乙出东门东行一百步，丙出南门南行五十一步而立。甲从乾隅东行三百二十步，望乙丙与城相参直。以太虚勾股测望。
}

\switchcolumn[0]*
\noindent\textbf{释曰}\quad 此以太虚勾股直接测望。乙东行为太虚勾，丙南行为太虚股。

\switchcolumn[1]
\noindent\textbf{【释】}\quad 这是用太虚勾股直接进行测望的问题。乙向东行走是太虚勾，丙向南行走是太虚股。

\switchcolumn[0]*
\noindent\textbf{术曰}\quad 太虚勾股各自乘相并开方得太虚弦。以通弦减之得城径。

\switchcolumn[1]
\noindent\textbf{【术】}\quad 太虚勾与太虚股各自乘，相加后开平方得到太虚弦。用通弦减去它得到城圆的直径。

\vspace{0.5em}

\Question{问八}{%
甲从乾隅东行三百二十步，南门外有槐树。乙出南门东至槐树下。甲斜望乙与槐与城相参直，计甲斜行一百零二步。求太虚股。
}

\switchcolumn[0]*
\noindent\textbf{释曰}\quad 此以太虚弦与明勾测望。甲斜行一百零二步为太虚弦，乙东至槐树下为明勾。

\switchcolumn[1]
\noindent\textbf{【释】}\quad 这是用太虚弦与明勾来进行测望的问题。甲斜着行走一百零二步是太虚弦，乙向东走到槐树下是明勾。

\switchcolumn[0]*
\noindent\textbf{术曰}\quad 太虚弦自乘，减明勾自乘，开方得太虚股。以太虚股减通股得城径。

\switchcolumn[1]
\noindent\textbf{【术】}\quad 太虚弦自乘，减去明勾自乘，开平方得到太虚股。用太虚股减通股得到城圆的直径。

\vspace{0.5em}

\Question{问九}{%
甲从艮隅南行过城门三百六十步而立，乙从坤隅东行过城门一百六十步而立。二人相望与城相参直。以太虚测望。
}

\switchcolumn[0]*
\noindent\textbf{释曰}\quad 此以小差股大差勾测望太虚。小差股大差勾即太虚勾股之属。

\switchcolumn[1]
\noindent\textbf{【释】}\quad 这是用小差股大差勾来测望太虚的问题。小差股大差勾就是太虚勾股这一类。

\switchcolumn[0]*
\noindent\textbf{术曰}\quad 小差股大差勾各自乘相并开方得太虚弦。以通弦减之得城径。

\switchcolumn[1]
\noindent\textbf{【术】}\quad 小差股与大差勾各自乘，相加后开平方得到太虚弦。用通弦减去它得到城圆的直径。

\vspace{0.5em}

\Question{问十}{%
乙出东门东行七十二步，丙出南门南行三十步而立。二人相望与城相参直。以太虚勾股测望。
}

\switchcolumn[0]*
\noindent\textbf{释曰}\quad 此以太虚勾股直接测望。乙东行为太虚勾，丙南行为太虚股。

\switchcolumn[1]
\noindent\textbf{【释】}\quad 这是用太虚勾股直接进行测望的问题。乙向东行走是太虚勾，丙向南行走是太虚股。

\switchcolumn[0]*
\noindent\textbf{术曰}\quad 太虚勾股各自乘相并开方得太虚弦。以勾股术求之得城径。

\switchcolumn[1]
\noindent\textbf{【术】}\quad 太虚勾与太虚股各自乘，相加后开平方得到太虚弦。用勾股术计算得到城圆的直径。

\vspace{0.5em}

\Question{问十一}{%
甲出东门东行二百步，乙出南门南行一百五十一步。甲望乙与城相参直，又斜行二百八十九步与乙相会。以太虚弦与皇极弦相较。
}

\switchcolumn[0]*
\noindent\textbf{释曰}\quad 此以太虚弦与皇极弦测望。甲斜行二百八十九步为皇极弦，太虚弦为其减数。

\switchcolumn[1]
\noindent\textbf{【释】}\quad 这是用太虚弦与皇极弦来进行测望的问题。甲斜着行走二百八十九步是皇极弦，太虚弦是比它小的数。

\switchcolumn[0]*
\noindent\textbf{术曰}\quad 皇极弦减太虚弦，余为明弦。以明弦减通弦得城径。

\switchcolumn[1]
\noindent\textbf{【术】}\quad 皇极弦减去太虚弦，余下就是明弦。用明弦减通弦得到城圆的直径。

\vspace{0.5em}

\Question{问十二}{%
甲从乾隅东行三百二十步，乙出东门南行三十步。甲望乙与城相参直。求太虚勾股较。
}

\switchcolumn[0]*
\noindent\textbf{释曰}\quad 此以通勾与太虚股测望。乙出东门南行为太虚股。

\switchcolumn[1]
\noindent\textbf{【释】}\quad 这是用通勾与太虚股来进行测望的问题。乙出东门向南行走是太虚股。

\switchcolumn[0]*
\noindent\textbf{术曰}\quad 太虚股自乘为实，以通勾约之得太虚勾。以减通勾得城径。

\switchcolumn[1]
\noindent\textbf{【术】}\quad 太虚股自乘作为被除数，用通勾来约得到太虚勾。用通勾减去它得到城圆的直径。

\vspace{0.5em}

\Question{问十三}{%
乙出南门东行七十二步，丙出东门南行三十步。甲从乾隅南行六百步，望乙丙与城相参直。以太虚测望。
}

\switchcolumn[0]*
\noindent\textbf{释曰}\quad 此以太虚勾股与通股测望。乙东行为太虚勾，丙南行为太虚股。

\switchcolumn[1]
\noindent\textbf{【释】}\quad 这是用太虚勾股与通股来进行测望的问题。乙向东行走是太虚勾，丙向南行走是太虚股。

\switchcolumn[0]*
\noindent\textbf{术曰}\quad 太虚勾股各自乘相并开方得太虚弦。以通弦减之得城径。

\switchcolumn[1]
\noindent\textbf{【术】}\quad 太虚勾与太虚股各自乘，相加后开平方得到太虚弦。用通弦减去它得到城圆的直径。

\vspace{0.5em}

\Question{问十四}{%
甲乙二人同出东门，甲东行乙南行。丙丁二人同出南门，丙南行丁东行。四人遥相望，计甲丙共行二百步，乙丁相距一百零二步。以太虚弦与皇极弦测望。
}

\switchcolumn[0]*
\noindent\textbf{释曰}\quad 此以太虚皇极二弦测望。乙丁相距为太虚弦，甲丙共行为皇极勾股和。

\switchcolumn[1]
\noindent\textbf{【释】}\quad 这是用太虚皇极二弦来进行测望的问题。乙丁相距是太虚弦，甲丙共行是皇极勾股和。

\switchcolumn[0]*
\noindent\textbf{术曰}\quad 太虚弦自乘，以皇极弦约之，得太虚勾。以减通勾得城径。

\switchcolumn[1]
\noindent\textbf{【术】}\quad 太虚弦自乘，用皇极弦来约，得到太虚勾。用通勾减去它得到城圆的直径。

%----------------------------------------------------------------------
\Section{明测望三}
%----------------------------------------------------------------------

\Question{问十五}{%
乙出南门东行七十二步而立。甲从乾隅南行六百步，望乙与城相参直。以明勾测望。
}

\switchcolumn[0]*
\noindent\textbf{释曰}\quad 此以明勾立法测望。乙出南门东行七十二步为明勾。

\switchcolumn[1]
\noindent\textbf{【释】}\quad 这是用明勾来建立方法进行测望的问题。乙出南门向东行走七十二步是明勾。

\switchcolumn[0]*
\noindent\textbf{术曰}\quad 明勾自乘为实，以通弦约之，得明股。以明股减通股得城径。

\switchcolumn[1]
\noindent\textbf{【术】}\quad 明勾自乘作为被除数，用通弦来约，得到明股。用明股减通股得到城圆的直径。

\vspace{0.5em}

\Question{问十六}{%
丙出南门南行五十一步而立。甲从乾隅东行三百二十步，望丙与城相参直。以明股测望。
}

\switchcolumn[0]*
\noindent\textbf{释曰}\quad 此以明股立法测望。丙出南门南行为明股。

\switchcolumn[1]
\noindent\textbf{【释】}\quad 这是用明股来建立方法进行测望的问题。丙出南门向南行走是明股。

\switchcolumn[0]*
\noindent\textbf{术曰}\quad 明股自乘为实，以通弦约之，得明勾。以明勾减通勾得城径。

\switchcolumn[1]
\noindent\textbf{【术】}\quad 明股自乘作为被除数，用通弦来约，得到明勾。用明勾减通勾得到城圆的直径。

\vspace{0.5em}

\Question{问十七}{%
乙出南门东行七十二步，丙出南门南行五十一步。二人相望与城相参直，计乙丙相距一百零二步。以明勾股测望。
}

\switchcolumn[0]*
\noindent\textbf{释曰}\quad 此以明勾股直接测望。乙东行为明勾，丙南行为明股，相距为明弦。

\switchcolumn[1]
\noindent\textbf{【释】}\quad 这是用明勾股直接进行测望的问题。乙向东行走是明勾，丙向南行走是明股，两人相距是明弦。

\switchcolumn[0]*
\noindent\textbf{术曰}\quad 明勾股各自乘相并开方得明弦。以通弦减之得城径。

\switchcolumn[1]
\noindent\textbf{【术】}\quad 明勾与明股各自乘，相加后开平方得到明弦。用通弦减去它得到城圆的直径。

\vspace{0.5em}

\Question{问十八}{%
甲从乾隅东行三百二十步，乙出南门东行七十二步。甲望乙与城相参直，又斜行二百八十九步与乙相会。以明勾与皇极弦测望。
}

\switchcolumn[0]*
\noindent\textbf{释曰}\quad 此以明勾与皇极弦测望。乙东行为明勾，甲斜行为皇极弦。

\switchcolumn[1]
\noindent\textbf{【释】}\quad 这是用明勾与皇极弦来进行测望的问题。乙向东行走是明勾，甲斜着行走是皇极弦。

\switchcolumn[0]*
\noindent\textbf{术曰}\quad 明勾自乘，减皇极弦自乘，开方得明股。以明股减通股得城径。

\switchcolumn[1]
\noindent\textbf{【术】}\quad 明勾自乘，减去皇极弦自乘，开平方得到明股。用明股减通股得到城圆的直径。

\vspace{0.5em}

\Question{问十九}{%
甲从乾隅南行六百步，丙出南门南行五十一步。甲望丙与城相参直。以明股与小差测望。
}

\switchcolumn[0]*
\noindent\textbf{释曰}\quad 此以明股与小差测望。丙南行为明股，小差即通弦减通勾之余。

\switchcolumn[1]
\noindent\textbf{【释】}\quad 这是用明股与小差来进行测望的问题。丙向南行走是明股，小差就是通弦减去通勾的余数。

\switchcolumn[0]*
\noindent\textbf{术曰}\quad 明股自乘，以小差约之，得明勾。以明勾减通勾得城径。

\switchcolumn[1]
\noindent\textbf{【术】}\quad 明股自乘，用小差来约，得到明勾。用明勾减通勾得到城圆的直径。

%----------------------------------------------------------------------
\Section{叀测望四}
%----------------------------------------------------------------------

\Question{问二十}{%
乙出东门南行三十步而立。甲从乾隅东行三百二十步，望乙与城相参直。以叀股测望。
}

\switchcolumn[0]*
\noindent\textbf{释曰}\quad 此以叀股立法测望。乙出东门南行为叀股。

\switchcolumn[1]
\noindent\textbf{【释】}\quad 这是用叀股来建立方法进行测望的问题。乙出东门向南行走是叀股。

\switchcolumn[0]*
\noindent\textbf{术曰}\quad 叀股自乘为实，以通弦约之，得叀勾。以叀勾减通勾得城径。

\switchcolumn[1]
\noindent\textbf{【术】}\quad 叀股自乘作为被除数，用通弦来约，得到叀勾。用叀勾减通勾得到城圆的直径。

\vspace{0.5em}

\Question{问二十一}{%
甲出东门东行二百步而立。丙出东门南行三十步。甲望丙与城相参直。以叀勾测望。
}

\switchcolumn[0]*
\noindent\textbf{释曰}\quad 此以叀勾立法测望。甲东行为叀勾，丙南行为叀股。

\switchcolumn[1]
\noindent\textbf{【释】}\quad 这是用叀勾来建立方法进行测望的问题。甲向东行走是叀勾，丙向南行走是叀股。

\switchcolumn[0]*
\noindent\textbf{术曰}\quad 叀勾自乘为实，以通弦约之，得叀股。以叀股减通股得城径。

\switchcolumn[1]
\noindent\textbf{【术】}\quad 叀勾自乘作为被除数，用通弦来约，得到叀股。用叀股减通股得到城圆的直径。

\vspace{0.5em}

\Question{问二十二}{%
甲出东门东行二百步，乙出东门南行三十步。二人相望与城相参直。以叀勾股直接测望。
}

\switchcolumn[0]*
\noindent\textbf{释曰}\quad 此以叀勾股直接测望。甲东行为叀勾，乙南行为叀股。

\switchcolumn[1]
\noindent\textbf{【释】}\quad 这是用叀勾股直接进行测望的问题。甲向东行走是叀勾，乙向南行走是叀股。

\switchcolumn[0]*
\noindent\textbf{术曰}\quad 叀勾股各自乘相并开方得叀弦。以通弦减之得城径。

\switchcolumn[1]
\noindent\textbf{【术】}\quad 叀勾与叀股各自乘，相加后开平方得到叀弦。用通弦减去它得到城圆的直径。

\vspace{0.5em}

\Question{问二十三}{%
甲从乾隅东行三百二十步，乙出东门南行三十步。甲斜行与乙相会，计甲斜行二百八十九步。以叀股与皇极弦测望。
}

\switchcolumn[0]*
\noindent\textbf{释曰}\quad 此以叀股与皇极弦测望。乙南行为叀股，甲斜行为皇极弦。

\switchcolumn[1]
\noindent\textbf{【释】}\quad 这是用叀股与皇极弦来进行测望的问题。乙向南行走是叀股，甲斜着行走是皇极弦。

\switchcolumn[0]*
\noindent\textbf{术曰}\quad 叀股自乘，减皇极弦自乘，开方得叀勾。以叀勾减通勾得城径。

\switchcolumn[1]
\noindent\textbf{【术】}\quad 叀股自乘，减去皇极弦自乘，开平方得到叀勾。用叀勾减通勾得到城圆的直径。

%----------------------------------------------------------------------
\Section{杂题测望五}
%----------------------------------------------------------------------

\Question{问二十四}{%
甲乙二人俱在城外西北乾隅，甲南行六百步而立，乙东行三百二十步而立。各不知余数。以勾股和较之变化，求杂题之数。
}

\switchcolumn[0]*
\noindent\textbf{释曰}\quad 此以杂题立法测望。杂题者，诸差互见之题也。

\switchcolumn[1]
\noindent\textbf{【释】}\quad 这是用杂题来建立方法进行测望的问题。杂题，就是各种差互相出现的题目。

\switchcolumn[0]*
\noindent\textbf{术曰}\quad 通勾股各自乘相并开方得通弦。以通弦减通股得大差，减通勾得小差。大小差相减得中差。诸差互以求之，得城径。

\switchcolumn[1]
\noindent\textbf{【术】}\quad 通勾与通股各自乘，相加后开平方得到通弦。用通弦减通股得到大差，减通勾得到小差。大小差相减得到中差。用各种差互相求解，得到城圆的直径。

\vspace{0.5em}

\Question{问二十五}{%
乙出东门南行三十步，丙出南门东行七十二步。二人相望与城相参直。以明叀二勾股测望。
}

\switchcolumn[0]*
\noindent\textbf{释曰}\quad 此以明叀二勾股杂测。乙南行为叀股，丙东行为明勾。

\switchcolumn[1]
\noindent\textbf{【释】}\quad 这是用明叀两种勾股进行杂测的问题。乙向南行走是叀股，丙向东行走是明勾。

\switchcolumn[0]*
\noindent\textbf{术曰}\quad 明勾叀股各自乘相并开方得杂弦。以通弦减之得城径。

\switchcolumn[1]
\noindent\textbf{【术】}\quad 明勾与叀股各自乘，相加后开平方得到杂弦。用通弦减去它得到城圆的直径。

\vspace{0.5em}

\Question{问二十六}{%
甲从乾隅东行三百二十步，乙出东门南行三十步，丙出南门东行七十二步。甲望乙丙与城相参直。以三差互见测望。
}

\switchcolumn[0]*
\noindent\textbf{释曰}\quad 此以大中小差杂见测望。乙南行为叀股，丙东行为明勾。

\switchcolumn[1]
\noindent\textbf{【释】}\quad 这是大中小差杂见进行测望的问题。乙向南行走是叀股，丙向东行走是明勾。

\switchcolumn[0]*
\noindent\textbf{术曰}\quad 明勾自乘为实，叀股自乘为实，二实相并开方得杂弦。以通弦减之得城径。

\switchcolumn[1]
\noindent\textbf{【术】}\quad 明勾自乘作为被除数，叀股自乘作为被除数，两个被除数相加后开平方得到杂弦。用通弦减去它得到城圆的直径。

\vspace{0.5em}

\Question{问二十七}{%
甲乙二人同出东门，甲东行乙南行。丙丁二人同出南门，丙南行丁东行。四人相望，计甲东行二百步，乙南行三十步，丙南行五十一步，丁东行七十二步。以四行互见测望。
}

\switchcolumn[0]*
\noindent\textbf{释曰}\quad 此以四行杂见测望。甲东行为叀勾，乙南行为叀股，丙南行为明股，丁东行为明勾。

\switchcolumn[1]
\noindent\textbf{【释】}\quad 这是用四行杂见来进行测望的问题。甲向东行走是叀勾，乙向南行走是叀股，丙向南行走是明股，丁向东行走是明勾。

\switchcolumn[0]*
\noindent\textbf{术曰}\quad 四行各自乘，相并开方得杂弦。以通弦减之得城径。

\switchcolumn[1]
\noindent\textbf{【术】}\quad 四种行走数各自乘，相加后开平方得到杂弦。用通弦减去它得到城圆的直径。

\vspace{0.5em}

\Question{问二十八}{%
甲从乾隅东行三百二十步，乙出东门南行三十步，丙出南门东行七十二步。甲斜望乙与城相参直，又斜望丙与城相参直。计甲至乙斜行二百八十九步，甲至丙斜行二百七十二步。以二斜行测望。
}

\switchcolumn[0]*
\noindent\textbf{释曰}\quad 此以二斜行杂测。甲至乙为皇极弦，甲至丙为小差弦。

\switchcolumn[1]
\noindent\textbf{【释】}\quad 这是用二斜行进行杂测的问题。甲到乙是皇极弦，甲到丙是小差弦。

\switchcolumn[0]*
\noindent\textbf{术曰}\quad 二弦各自乘相减，开方得勾股较。以加减术推之得城径。

\switchcolumn[1]
\noindent\textbf{【术】}\quad 两条弦各自乘，相减后开平方得到勾股差。用加减术推算得到城圆的直径。

\vspace{0.5em}

\Question{问二十九}{%
乙出东门东行二百步，丙出南门南行五十一步。二人相望与城相参直。甲从乾隅东行三百二十步，望乙丙与城相参直。以东行南行杂测。
}

\switchcolumn[0]*
\noindent\textbf{释曰}\quad 此以东行南行杂见测望。乙东行为边勾，丙南行为明股。

\switchcolumn[1]
\noindent\textbf{【释】}\quad 这是用东行南行杂见来进行测望的问题。乙向东行走是边勾，丙向南行走是明股。

\switchcolumn[0]*
\noindent\textbf{术曰}\quad 边勾明股各自乘相并开方得杂弦。以通弦减之得城径。

\switchcolumn[1]
\noindent\textbf{【术】}\quad 边勾与明股各自乘，相加后开平方得到杂弦。用通弦减去它得到城圆的直径。

\vspace{0.5em}

\Question{问三十}{%
甲从艮隅南行过城门三百六十步，乙从坤隅东行过城门一百六十步。二人相望与城相参直。以过城门之行杂测。
}

\switchcolumn[0]*
\noindent\textbf{释曰}\quad 此以过城门之行杂测。甲南行过城门为小差股，乙东行过城门为大差勾。

\switchcolumn[1]
\noindent\textbf{【释】}\quad 这是用走过城门的行走数进行杂测的问题。甲向南走过城门是小差股，乙向东走过城门是大差勾。

\switchcolumn[0]*
\noindent\textbf{术曰}\quad 大差勾小差股各自乘相并开方得杂弦。以通弦减之得城径。

\switchcolumn[1]
\noindent\textbf{【术】}\quad 大差勾与小差股各自乘，相加后开平方得到杂弦。用通弦减去它得到城圆的直径。

\vspace{0.5em}

\Question{问三十一}{%
甲乙二人俱在城外西北乾隅，甲南行六百步而立，乙东行三百二十步而立。甲斜行与乙相会。以通勾股弦诸差互求。
}

\switchcolumn[0]*
\noindent\textbf{释曰}\quad 此以诸差互见杂测。通勾股弦诸差互相求之術。

\switchcolumn[1]
\noindent\textbf{【释】}\quad 这是用诸差互见进行杂测的问题。通勾股弦各种差互相求解的方法。

\switchcolumn[0]*
\noindent\textbf{术曰}\quad 通弦自乘，减通股自乘，开方得大差勾。通弦自乘，减通勾自乘，开方得小差股。大小差相减得中差。诸差互以求之，得城径。

\switchcolumn[1]
\noindent\textbf{【术】}\quad 通弦自乘，减去通股自乘，开平方得到大差勾。通弦自乘，减去通勾自乘，开平方得到小差股。大小差相减得到中差。用各种差互相求解，得到城圆的直径。

\vspace{0.5em}

\Question{问三十二}{%
甲出东门东行二百步，乙出南门南行五十一步。二人相望与城相参直，计相距二百八十九步。以相距与各行杂测。
}

\switchcolumn[0]*
\noindent\textbf{释曰}\quad 此以相距与各行杂见测望。相距为皇极弦。

\switchcolumn[1]
\noindent\textbf{【释】}\quad 这是用相距与各行杂见来进行测望的问题。相距就是皇极弦。

\switchcolumn[0]*
\noindent\textbf{术曰}\quad 相距自乘，减东行自乘，开方得南行余数。以减通股得城径。

\switchcolumn[1]
\noindent\textbf{【术】}\quad 相距自乘，减去东行自乘，开平方得到南行的余数。用通股减去它得到城圆的直径。

\vspace{0.5em}

\Question{问三十三}{%
乙出东门南行三十步，丙出南门东行七十二步，丁出东门东行二百步。甲从乾隅东行三百二十步，望乙丙丁与城相参直。以三行杂测。
}

\switchcolumn[0]*
\noindent\textbf{释曰}\quad 此以三行杂见测望。乙南行为叀股，丙东行为明勾，丁东行为边勾。

\switchcolumn[1]
\noindent\textbf{【释】}\quad 这是用三行杂见来进行测望的问题。乙向南行走是叀股，丙向东行走是明勾，丁向东行走是边勾。

\switchcolumn[0]*
\noindent\textbf{术曰}\quad 三行各自乘，相并开方得杂弦。以通弦减之得城径。

\switchcolumn[1]
\noindent\textbf{【术】}\quad 三种行走数各自乘，相加后开平方得到杂弦。用通弦减去它得到城圆的直径。

\vspace{0.5em}

\Question{问三十四}{%
甲乙丙丁四人，甲从乾隅东行三百二十步，乙从乾隅南行六百步，丙出东门东行二百步，丁出南门南行五十一步。以四隅之行杂测。
}

\switchcolumn[0]*
\noindent\textbf{释曰}\quad 此以四隅之行杂见测望。诸行互见，以天元术推之。

\switchcolumn[1]
\noindent\textbf{【释】}\quad 这是用四隅之行杂见来进行测望的问题。各种行走数互相出现，用天元术来推算。

\switchcolumn[0]*
\noindent\textbf{术曰}\quad 通勾股各自乘相并开方得通弦。以通弦减通股得大差，减通勾得小差。以大小差推诸数，得城径。

\switchcolumn[1]
\noindent\textbf{【术】}\quad 通勾与通股各自乘，相加后开平方得到通弦。用通弦减通股得到大差，减通勾得到小差。用大小差推算各种数值，得到城圆的直径。

\EndVol{测圆海镜分类释术卷八终}

%======================================================================
\Volume{九}
%======================================================================

\switchcolumn[0]*
\noindent\textbf{测圆海镜分类释术卷九}\quad 元~李~冶~撰\\
明~顾应祥~释术

\switchcolumn[1]
\noindent\textbf{测圆海镜分类释术·第九卷}\quad 元代~李~冶~撰\\
明代~顾应祥~释术

\vspace{1em}

\switchcolumn[0]*
卷九专论诸弦差较之测望。诸弦者，通弦、边弦、底弦、黄广弦、黄长弦之类；差较者，弦与勾股之差及差与差之比较也。凡十八题，分为五节。

\switchcolumn[1]
第九卷专门论述各种弦差差的测望方法。诸弦，包括通弦、边弦、底弦、黄广弦、黄长弦等；差较，就是弦与勾股的差以及差与差的比较。共有十八道题，分为五节。

%----------------------------------------------------------------------
\Section{大股小勾测望一}
%----------------------------------------------------------------------

\Question{问一}{%
甲从乾隅南行六百步而立，乙从乾隅东行三百二十步而立。以通股与通勾之差，求大股小勾。
}

\switchcolumn[0]*
\noindent\textbf{释曰}\quad 此以大股小勾立法测望。通股六百步为大股，通勾三百二十步为小勾。大股小勾之相较，为勾股之极差。

\switchcolumn[1]
\noindent\textbf{【释】}\quad 这是用大股小勾来建立方法进行测望的问题。通股六百步是大股，通勾三百二十步是小勾。大股与小勾的比较，就是勾股的极差。

\switchcolumn[0]*
\noindent\textbf{术曰}\quad 通股自乘，减通勾自乘，开方得勾股和。以和减通弦，余为弦较。以弦较加减通勾股，得城径。

\switchcolumn[1]
\noindent\textbf{【术】}\quad 通股自乘，减去通勾自乘，开平方得到勾股和。用和减通弦，余下就是弦较。用弦较加减通勾股，得到城圆的直径。

\vspace{0.5em}

\Question{问二}{%
乙出东门南行三十步而立，丙出南门东行七十二步而立。甲从乾隅东行三百二十步，望乙丙与城相参直。以大股小勾测望。
}

\switchcolumn[0]*
\noindent\textbf{释曰}\quad 此以叀股明勾测望大股小勾。乙南行为叀股，丙东行为明勾。

\switchcolumn[1]
\noindent\textbf{【释】}\quad 这是用叀股明勾来测望大股小勾的问题。乙向南行走是叀股，丙向东行走是明勾。

\switchcolumn[0]*
\noindent\textbf{术曰}\quad 叀股明勾各自乘相并开方得小弦。以通弦减之得城径。

\switchcolumn[1]
\noindent\textbf{【术】}\quad 叀股与明勾各自乘，相加后开平方得到小弦。用通弦减去它得到城圆的直径。

\vspace{0.5em}

\Question{问三}{%
甲从乾隅东行三百二十步，乙出东门东行二百步。二人相望与城相参直。以大股与边勾测望。
}

\switchcolumn[0]*
\noindent\textbf{释曰}\quad 此以大股边勾测望。甲东行为通勾，乙东行为边勾。

\switchcolumn[1]
\noindent\textbf{【释】}\quad 这是用大股边勾来进行测望的问题。甲向东行走是通勾，乙向东行走是边勾。

\switchcolumn[0]*
\noindent\textbf{术曰}\quad 通勾减边勾，余为小勾。以小勾加减术推之得城径。

\switchcolumn[1]
\noindent\textbf{【术】}\quad 通勾减去边勾，余下就是小勾。用小勾加减术推算得到城圆的直径。

%----------------------------------------------------------------------
\Section{大差大小差测望二}
%----------------------------------------------------------------------

\Question{问四}{%
甲乙二人俱在城外西北乾隅，甲南行六百步而立，乙东行三百二十步而立。以通弦减通股为大差，减通勾为小差。大小差相求。
}

\switchcolumn[0]*
\noindent\textbf{释曰}\quad 此以大差大小差立法测望。大差即通股弦较，小差即通勾弦较。

\switchcolumn[1]
\noindent\textbf{【释】}\quad 这是用大差大小差来建立方法进行测望的问题。大差就是通股弦差，小差就是通勾弦差。

\switchcolumn[0]*
\noindent\textbf{术曰}\quad 大差自乘为实，小差自乘为实，二实相减开方得中差。以中差加减大小差，得城径。

\switchcolumn[1]
\noindent\textbf{【术】}\quad 大差自乘作为被除数，小差自乘作为被除数，两个被除数相减后开平方得到中差。用中差加减大小差，得到城圆的直径。

\vspace{0.5em}

\Question{问五}{%
乙出东门南行三十步而立，丙出南门东行七十二步而立。以叀股明勾求大小差。
}

\switchcolumn[0]*
\noindent\textbf{释曰}\quad 此以叀股明勾测望大小差。叀股即小差之属，明勾即大差之属。

\switchcolumn[1]
\noindent\textbf{【释】}\quad 这是用叀股明勾来测望大小差的问题。叀股就是小差这一类，明勾就是大差这一类。

\switchcolumn[0]*
\noindent\textbf{术曰}\quad 叀股自乘减明勾自乘，开方得大小差较。以较加减术推之得城径。

\switchcolumn[1]
\noindent\textbf{【术】}\quad 叀股自乘减去明勾自乘，开平方得到大小差差。用差加减术推算得到城圆的直径。

\vspace{0.5em}

\Question{问六}{%
甲从乾隅东行三百二十步，乙出东门南行三十步，丙出南门东行七十二步。甲望乙丙与城相参直。以三差互求。
}

\switchcolumn[0]*
\noindent\textbf{释曰}\quad 此以大中小差互见测望。乙南行为叀股，丙东行为明勾。

\switchcolumn[1]
\noindent\textbf{【释】}\quad 这是大中小差互见进行测望的问题。乙向南行走是叀股，丙向东行走是明勾。

\switchcolumn[0]*
\noindent\textbf{术曰}\quad 明勾自乘为大差实，叀股自乘为小差实。二实相减开方得中差。以中差加减得城径。

\switchcolumn[1]
\noindent\textbf{【术】}\quad 明勾自乘作为大差被除数，叀股自乘作为小差被除数。两个被除数相减后开平方得到中差。用中差加减得到城圆的直径。

\vspace{0.5em}

\Question{问七}{%
甲从艮隅南行过城门三百六十步，乙从坤隅东行过城门一百六十步。以过城门之行求大小差。
}

\switchcolumn[0]*
\noindent\textbf{释曰}\quad 此以过城门之行测望大小差。甲南行为小差股，乙东行为大差勾。

\switchcolumn[1]
\noindent\textbf{【释】}\quad 这是用走过城门的行走数来测望大小差的问题。甲向南行走是小差股，乙向东行走是大差勾。

\switchcolumn[0]*
\noindent\textbf{术曰}\quad 小差股自乘为实，大差勾自乘为实。二实相并开方得大小差弦。以通弦减之得城径。

\switchcolumn[1]
\noindent\textbf{【术】}\quad 小差股自乘作为被除数，大差勾自乘作为被除数。两个被除数相加后开平方得到大小差弦。用通弦减去它得到城圆的直径。

%----------------------------------------------------------------------
\Section{诸弦测望三}
%----------------------------------------------------------------------

\Question{问八}{%
甲乙二人俱在城外西北乾隅，甲南行六百步而立，乙东行三百二十步而立。甲斜行与乙相会。以通弦为本，求边弦、底弦。
}

\switchcolumn[0]*
\noindent\textbf{释曰}\quad 此以诸弦立法测望。通弦为勾股弦之全，边弦、底弦为其分支。

\switchcolumn[1]
\noindent\textbf{【释】}\quad 这是用诸弦来建立方法进行测望的问题。通弦是勾股弦的全体，边弦、底弦是它的分支。

\switchcolumn[0]*
\noindent\textbf{术曰}\quad 通勾股各自乘相并开方得通弦。通弦半之得皇极弦。以通弦减皇极弦余为边弦。边弦减通弦余为城径。

\switchcolumn[1]
\noindent\textbf{【术】}\quad 通勾与通股各自乘，相加后开平方得到通弦。通弦取一半得到皇极弦。用通弦减皇极弦余下就是边弦。边弦减通弦余下就是城圆的直径。

\vspace{0.5em}

\Question{问九}{%
乙出东门东行二百步，丙出南门南行五十一步。二人相望与城相参直。以边股明勾求诸弦。
}

\switchcolumn[0]*
\noindent\textbf{释曰}\quad 此以边股明勾测望诸弦。乙东行为边股，丙南行为明股。

\switchcolumn[1]
\noindent\textbf{【释】}\quad 这是用边股明勾来测望诸弦的问题。乙向东行走是边股，丙向南行走是明股。

\switchcolumn[0]*
\noindent\textbf{术曰}\quad 边股明股各自乘相并开方得边弦。以通弦减之得城径。

\switchcolumn[1]
\noindent\textbf{【术】}\quad 边股与明股各自乘，相加后开平方得到边弦。用通弦减去它得到城圆的直径。

\vspace{0.5em}

\Question{问十}{%
甲从乾隅东行三百二十步，乙出东门南行三十步。甲望乙与城相参直，又斜行二百八十九步与乙相会。以底弦与通弦测望。
}

\switchcolumn[0]*
\noindent\textbf{释曰}\quad 此以底弦通弦测望。甲斜行为皇极弦，即底弦之属。

\switchcolumn[1]
\noindent\textbf{【释】}\quad 这是用底弦通弦来进行测望的问题。甲斜着行走是皇极弦，就是底弦这一类。

\switchcolumn[0]*
\noindent\textbf{术曰}\quad 皇极弦自乘，减通勾自乘，开方得底股。以底股减通股得城径。

\switchcolumn[1]
\noindent\textbf{【术】}\quad 皇极弦自乘，减去通勾自乘，开平方得到底股。用底股减通股得到城圆的直径。

%----------------------------------------------------------------------
\Section{大小差测望四}
%----------------------------------------------------------------------

\Question{问十一}{%
甲乙二人俱在城外西北乾隅，甲南行六百步而立，乙东行三百二十步而立。以通股减通弦为大差，通勾减通弦为小差。大小差互求。
}

\switchcolumn[0]*
\noindent\textbf{释曰}\quad 此以大小差互见测望。大差为通股弦较，小差为通勾弦较。

\switchcolumn[1]
\noindent\textbf{【释】}\quad 这是用大小差互见来进行测望的问题。大差是通股弦差，小差是通勾弦差。

\switchcolumn[0]*
\noindent\textbf{术曰}\quad 大差小差各自乘相并开方得中差弦。以通弦减之得城径。

\switchcolumn[1]
\noindent\textbf{【术】}\quad 大差与小差各自乘，相加后开平方得到中差弦。用通弦减去它得到城圆的直径。

\vspace{0.5em}

\Question{问十二}{%
乙出东门南行三十步，丙出南门东行七十二步。甲从乾隅东行三百二十步，望乙丙与城相参直。以叀股明勾求大小差。
}

\switchcolumn[0]*
\noindent\textbf{释曰}\quad 此以叀股明勾测望大小差。

\switchcolumn[1]
\noindent\textbf{【释】}\quad 这是用叀股明勾来测望大小差的问题。

\switchcolumn[0]*
\noindent\textbf{术曰}\quad 叀股自乘为小差实，明勾自乘为大差实。二实相减开方得差较。以差较加减得城径。

\switchcolumn[1]
\noindent\textbf{【术】}\quad 叀股自乘作为小差被除数，明勾自乘作为大差被除数。两个被除数相减后开平方得到差较。用差较加减得到城圆的直径。

\vspace{0.5em}

\Question{问十三}{%
甲从艮隅南行过城门三百六十步，乙从坤隅东行过城门一百六十步。以过城门之行求差较。
}

\switchcolumn[0]*
\noindent\textbf{释曰}\quad 此以过城门之行测望差较。

\switchcolumn[1]
\noindent\textbf{【释】}\quad 这是用走过城门的行走数来测望差较的问题。

\switchcolumn[0]*
\noindent\textbf{术曰}\quad 小差股大差勾各自乘相并开方得差弦。以通弦减之得城径。

\switchcolumn[1]
\noindent\textbf{【术】}\quad 小差股与大差勾各自乘，相加后开平方得到差弦。用通弦减去它得到城圆的直径。

\vspace{0.5em}

\Question{问十四}{%
甲出东门东行二百步，乙出南门南行五十一步。二人相望与城相参直。以边股明股求大小差。
}

\switchcolumn[0]*
\noindent\textbf{释曰}\quad 此以边股明股测望大小差。

\switchcolumn[1]
\noindent\textbf{【释】}\quad 这是用边股明股来测望大小差的问题。

\switchcolumn[0]*
\noindent\textbf{术曰}\quad 边股明股各自乘相并开方得弦。以通弦减之得城径。

\switchcolumn[1]
\noindent\textbf{【术】}\quad 边股与明股各自乘，相加后开平方得到弦。用通弦减去它得到城圆的直径。

\EndVol{测圆海镜分类释术卷九终}

%======================================================================
\Volume{十}
%======================================================================

\switchcolumn[0]*
\noindent\textbf{测圆海镜分类释术卷十}\quad 元~李~冶~撰\\
明~顾应祥~释术

\switchcolumn[1]
\noindent\textbf{测圆海镜分类释术·第十卷}\quad 元代~李~冶~撰\\
明代~顾应祥~释术

\vspace{1em}

\switchcolumn[0]*
卷十为全书之终，专论杂法。杂法者，诸术之变通也，凡不属于前九卷之专题者，皆汇于此。凡四题，皆为通术之总结。

\switchcolumn[1]
第十卷是全书的终结，专门论述杂法。杂法，就是各种方法的变通，凡是不属于前九卷专题的内容，都汇集在这里。共有四道题，都是通用方法的总结。

%----------------------------------------------------------------------
\Section{杂法}
%----------------------------------------------------------------------

\Question{问一}{%
甲乙二人俱在城外西北乾隅，甲南行六百步而立，乙东行三百二十步而立。二人隔城相望，甲斜行与乙相会。以天元术总求诸数。
}

\switchcolumn[0]*
\noindent\textbf{释曰}\quad 此以天元术杂法总测。天元一者，立天元一为勾股之基，以诸数推之。

\switchcolumn[1]
\noindent\textbf{【释】}\quad 这是用天元术杂法进行总体测望的问题。天元一，就是设天元一作为勾股的基础，用各个数值来推算。

\switchcolumn[0]*
\noindent\textbf{术曰}\quad 立天元一为勾，以股乘之，得勾股积。以弦除之，得容圆径。天元自之为勾幂，以减弦幂，开方得股。勾股各自乘相并开方得弦。以勾股求容圆术，得城径二百四十步。

\switchcolumn[1]
\noindent\textbf{【术】}\quad 设天元一作为勾，用股来乘，得到勾股积。用弦来除，得到容圆的直径。天元自乘作为勾的幂，用来减弦的幂，开平方得到股。勾与股各自乘，相加后开平方得到弦。用勾股求容圆的方法，得到城圆直径二百四十步。

\vspace{0.5em}

\Question{问二}{%
乙出东门南行三十步，丙出南门东行七十二步，丁出东门东行二百步，戊出南门南行五十一步。五人各不知余数。甲从乾隅东行三百二十步，望乙丙丁戊与城相参直。以天元术杂测诸行。
}

\switchcolumn[0]*
\noindent\textbf{释曰}\quad 此以天元术杂测诸行。乙南行为叀股，丙东行为明勾，丁东行为边勾，戊南行为明股。

\switchcolumn[1]
\noindent\textbf{【释】}\quad 这是用天元术来杂测各种行走数的问题。乙向南行走是叀股，丙向东行走是明勾，丁向东行走是边勾，戊向南行走是明股。

\switchcolumn[0]*
\noindent\textbf{术曰}\quad 立天元一为城径。以通勾减天元，得边勾。以通股减天元，得明股。边勾明股各自乘相并开方得皇极弦。以通弦减皇极弦，余为太虚弦。太虚弦自乘，以明和约之，得太虚勾。以次推之，得天元即城径二百四十步。

\switchcolumn[1]
\noindent\textbf{【术】}\quad 设天元一作为城圆的直径。用通勾减天元，得到边勾。用通股减天元，得到明股。边勾与明股各自乘，相加后开平方得到皇极弦。用通弦减皇极弦，余下就是太虚弦。太虚弦自乘，用明和来约，得到太虚勾。依次推算，得到天元就是城圆直径二百四十步。

\vspace{0.5em}

\Question{问三}{%
甲乙丙丁四人，甲从乾隅东行三百二十步，乙从乾隅南行六百步，丙出东门东行二百步，丁出南门南行五十一步。以天元术总括全书之术。
}

\switchcolumn[0]*
\noindent\textbf{释曰}\quad 此以天元术总括全书。通勾三百二十步，通股六百步，通弦以勾股术求之。边勾、明股、皇极弦、太虚弦，皆以天元一推之。

\switchcolumn[1]
\noindent\textbf{【释】}\quad 这是用天元术来总括全书的问题。通勾三百二十步，通股六百步，通弦用勾股术来求。边勾、明股、皇极弦、太虚弦，都用天元一来推算。

\switchcolumn[0]*
\noindent\textbf{术曰}\quad 立天元一为圆径。通勾自乘，以天元减之，得边勾幂。通股自乘，以天元减之，得明股幂。边勾幂明股幂各自乘相并开方得皇极弦。皇极弦自乘，以通弦减之，得诸弦之差。以次推之，得天元即城径二百四十步。

\switchcolumn[1]
\noindent\textbf{【术】}\quad 设天元一作为圆的直径。通勾自乘，用天元来减，得到边勾的幂。通股自乘，用天元来减，得到明股的幂。边勾幂与明股幂各自乘，相加后开平方得到皇极弦。皇极弦自乘，用通弦来减，得到各种弦的差。依次推算，得到天元就是城圆直径二百四十步。

\vspace{0.5em}

\Question{问四}{%
全书诸题，以勾股容圆为本，以天元一为法。通勾、通股、通弦为三大纲，边勾、明股、皇极弦、太虚弦为四目，大差、小差、中差为三变。总括其术。
}

\switchcolumn[0]*
\noindent\textbf{释曰}\quad 此总括全书之术。测圆海镜分类释术，凡十卷，一百七十题。以勾股容圆为题，以天元术为法。通勾三百二十步，通股六百步，通弦以勾股术求之得六百八十步。城径二百四十步。诸差：大差八十步，小差三百六十步，中差二百八十步。边勾一百六十步，明股七十二步，皇极弦二百八十九步，太虚弦一百零二步。凡此诸数，皆以天元一推之而得。

\switchcolumn[1]
\noindent\textbf{【释】}\quad 这是总括全书方法的问题。测圆海镜分类释术，共十卷，一百七十道题。以勾股容圆为题目，以天元术为方法。通勾三百二十步，通股六百步，通弦用勾股术求得六百八十步。城圆直径二百四十步。各种差：大差八十步，小差三百六十步，中差二百八十步。边勾一百六十步，明股七十二步，皇极弦二百八十九步，太虚弦一百零二步。所有这些数值，都是用天元一推算而得的。

\switchcolumn[0]*
\noindent\textbf{术曰}\quad 立天元一为圆径。以勾股容圆术推之：勾股和减弦，余为圆径。通勾三百二十步，通股六百步，通弦六百八十步。勾股和九百二十步，减弦六百八十步，余二百四十步，即城径。此全书之宗也。

\switchcolumn[1]
\noindent\textbf{【术】}\quad 设天元一作为圆的直径。用勾股容圆术来推算：勾股和减去弦，余下就是圆的直径。通勾三百二十步，通股六百步，通弦六百八十步。勾股和九百二十步，减去弦六百八十步，余下二百四十步，就是城圆的直径。这是全书的宗旨。

\EndVol{测圆海镜分类释术卷十终}

\end{paracol}

\vspace{2em}
\begin{center}
\rule{0.6\textwidth}{0.5pt}

\textbf{测圆海镜分类释术卷六至卷十终}

\vspace{0.5em}

{\small 元李冶撰\quad 明顾应祥释术}

\vspace{0.5em}

{\small 钦定四库全书\quad 子部\quad 天文算法类}
\end{center}

\end{document}
